\section{The maximal comparison and its scope}
\label{four:maximal:section}

Maximality allows the boundary deformation to retain a decaying second
fundamental form all the way to infinity.  The reason is the homogeneous
trace equation in Section~\ref{b:transition}: when $\tr_gK=0$, the tensor
contributes a drift proportional to the graph gradient.  Thus the graph
can decay exponentially even though $K$ does not vanish outside a compact
set.  The conformal equation still sees $|K|^2/2$.  We now construct the
solutions for which that end calculation applies.

There are two geometric tasks before solving the equations.  We must
create strict dominant-energy slack without destroying maximality or
changing the limiting energy and enclosing area.  We must then select an
inner boundary whose outward collars control the large graph slopes.
The remaining analysis has the same input and output as the boundary
method: nonnegative scalar curvature, negative inner mean curvature,
a lower metric comparison on every cut, and an upper error in energy.
The noncompact drift requires the tail weights and continuation estimates
proved in this part.

The following theorem records the maximal-vacuum application considered
here.  Corollary~\ref{b:four-weak} gives a broader numerical comparison
under its one-ended weak-decay and dominant-energy hypotheses.  That
corollary first replaces the distant end and tensor.  The construction
below retains the decaying tensor tail through the deformation.

Throughout, four refers to the spatial dimension. The associated spacetime
dimension is five, and a horizon has dimension three. Write
\(\omega_3=|S^3|=2\pi^2\). The notation \(O_j(r^{-q})\) includes coordinate
derivatives through order \(j\), with the corresponding additional decay.

\begin{theorem}\label{four:maximal:thm:main}
Let \((M,g,K)\) be a smooth connected orientable initial-data exterior,
diffeomorphic to \(\R^4\) minus an open ball, complete up to its connected
compact boundary \(S\cong S^3\), and with one asymptotically flat end.
Here \(K\) is a symmetric covariant two-tensor. Assume the vacuum constraint
equations with zero cosmological constant and maximality:
\begin{equation}\label{four:maximal:intro:constraints}
 R_g+(\tr_gK)^2-|K|_g^2=0,\qquad
 \divg_g\bigl(K-(\tr_gK)g\bigr)=0,\qquad \tr_gK=0.
\end{equation}
Assume a smooth effective \(T^2=U(1)\times U(1)\) action preserving \(g,K\)
which, in an asymptotically Euclidean chart \(\R^4=\mathbb C^2\), is
\[
 (e^{i\alpha},e^{i\beta})\cdot(z_1,z_2)
       =(e^{i\alpha}z_1,e^{i\beta}z_2)
\]
at all sufficiently large radii. No cohomogeneity-one assumption is made.
In that chart suppose
\begin{equation}\label{four:maximal:intro:decay}
 g_{ij}-\delta_{ij}=O_2(r^{-q}),\qquad
 K_{ij}=O_1(r^{-q-1}),\qquad q>1.
\end{equation}
Suppose the following ADM limits exist and are finite:
\begin{align}
 E&=\frac{1}{6\omega_3}\lim_{r\to\infty}
     \int_{|x|=r}(\partial_jg_{ij}-\partial_ig_{jj})
                  \nu_\delta^i\dd A_\delta,\label{four:maximal:intro:energy}\\
 P_i&=\frac{1}{3\omega_3}\lim_{r\to\infty}
     \int_{|x|=r}\bigl(K_{ij}-(\tr_gK)g_{ij}\bigr)
                  \nu_\delta^j\dd A_\delta.\label{four:maximal:intro:momentum}
\end{align}
Assume \(P=0\) and \(E>0\).

Orient \(S\) by its unit normal \(\nu\) toward the end and set
\(H_S=\divg_S\nu\). Suppose
\begin{equation}\label{four:maximal:intro:mots}
 \Theta_K(S):=H_S+\tr_SK=0.
\end{equation}
Suppose \(S\) is outermost toward the end: there is no compact smooth
embedded two-sided hypersurface in the interior enclosing \(S\), oriented
toward the end, with \(\Theta_K\le0\) everywhere. Suppose also that \(S\)
is outer area-minimizing: every compact smooth embedded enclosing
hypersurface has three-volume at least \(A=|S|_g\). Disconnected and
noninvariant enclosing competitors are included.
Then
\begin{equation}\label{four:maximal:intro:main-inequality}
 E\ge\frac12\left(\frac{A}{2\pi^2}\right)^{2/3}.
\end{equation}
\end{theorem}

An enclosing hypersurface separates \(S\) from the selected end and bounds
with \(S\) a compact region. The original \(S\) is itself allowed when
taking the enclosing-area infimum. All hypersurface areas in this part
are three-dimensional volumes. In particular, outer area-minimization in
Theorem~\ref{four:maximal:thm:main} is not restricted by the torus action.
The torus action and outermostness specify the stated application class;
the construction below does not use them.  Ball-exterior topology supplies
the compact filling, and vacuum supplies the initial strictification.
Outer area-minimization identifies the final enclosing-area infimum with
\(A\).  Neither \(H_S=0\) nor \(\tr_SK=0\) is imposed.  The conditions
\(P=0\) and \(E>0\) identify the energy in
\eqref{four:maximal:intro:energy} with the ADM rest mass.

The normalization is sharp. The time-symmetric Schwarzschild--Tangherlini
family \cite{Tangherlini1963} has, in isotropic coordinates of radius
\(\varrho\), the spatial metric
\[
 g_m=\left(1+\frac{m}{2\varrho^2}\right)^2\delta,
 \qquad K=0,
 \qquad \varrho\ge\sqrt{m/2};
\]
see also \cite[p.~83]{BrayLee2009}. Its leading metric term is
\(m\varrho^{-2}\delta\), so \eqref{four:maximal:intro:energy} gives energy \(m\).
The boundary sphere has volume
\(8\omega_3(m/2)^{3/2}=\omega_3(2m)^{3/2}\), making the right side of
\eqref{four:maximal:intro:main-inequality} equal to \(m\).


Unless a different metric is indicated, norms, contractions, covariant
derivatives and index raising use the background metric \(g\). We use
\(\Delta=\divg\grad\). For symmetric two-tensors define
\begin{align}
 P_D&=D-(\tr_gD)g, &
 \mathsf{B}(D_1,D_2)&=D_1:D_2-(\tr_gD_1)(\tr_gD_2),\label{four:maximal:intro:bilinear}\\
 \mu&=\tfrac12\bigl(R_g-\mathsf{B}(K,K)\bigr), &
 J&=\divg_gP_K.\label{four:maximal:intro:densities}
\end{align}
Thus the vacuum constraints give \(\mu=J=0\).
For an oriented hypersurface the expansion is
\(\Theta_K=H+\tr_{\mathrm{tan}}K\), with the normal toward the
specified exterior. On an inner boundary this normal points into the
exterior domain; on a distant coordinate sphere the outward normal points
toward infinity. We retain these conventions in every boundary flux.

For an inner cut \(S_0\), let
\[
 \mathcal A_g(S_0)=\inf\{|\Sigma|_g:
              \Sigma\text{ is a smooth enclosing cut of }S_0\},
\]
using the full class of cuts in Definition~\ref{four:def:cuts} and
including \(S_0\) when it is smooth. The proof uses this infimum until
the final application of the outer area-minimizing hypothesis.


\subsection{The geometric output and the route to it}

It is useful to state what the construction will produce for strict
data before introducing its equations.  This is the point at which
the analytic problem rejoins the mass--area comparison.

\begin{proposition}[Maximal deformation output]
\label{four:maximal:prop:geometric-output}
Let $(M,g,K)$ be a smooth connected orientable four-dimensional
ball exterior, complete including its compact boundary $S$, with one
Euclidean coordinate end.  Let $r\ge1$ be a fixed smooth positive
extension of its Euclidean radius.  Suppose $1<q<2$, $0<\delta<q-1$,
\[
 \tr_gK=0,\quad g-\delta_{\rm Eucl}=O_2(r^{-q}),\quad
 K=O_1(r^{-1-q}),\quad R_g=O(r^{-4-\delta}),
 \qquad \mu-|J|\ge c_0r^{-4-\delta}>0,
\]
and suppose $H_S+\tr_SK<0$.  Assume finite ADM energy and uniform
bounds of every fixed derivative order on end patches of unit size.
Write $A_* =\mathcal A_g(S)$ for the infimum over all smooth enclosing
cuts, including disconnected cuts.

There is a smooth cut $B$ enclosing $S$, with connected one-ended
exterior $\Omega$, such that for each $0<\eps<1$ and all sufficiently
large integers $N$ there is a smooth metric $\widehat g_N$ on
$\overline\Omega$ satisfying
\[
 \widehat g_N\ge e^{-2\eps}g,\qquad
 R_{\widehat g_N}\ge0,\qquad \widehat H_B<0,
 \qquad
 E_{\widehat g_N}\le E_g+
       \frac{\Pi_N}{3\omega_3}
            (e^{-\eps N}+e^{-\eps N/2}).
\]
Here $\Pi_N$ has polynomial growth at fixed strict data and $\eps$.
The new metric is complete including $B$, has
$\widehat g_N-\delta_{\rm Eucl}=O_2(r^{-q})$ and
$R_{\widehat g_N}=O(r^{-4-\delta})$, and every smooth cut enclosing
$B$ has $\widehat g_N$-area at least $e^{-3\eps}A_*$.
\end{proposition}

The proof occupies Sections~\ref{four:maximal:sec:preparation}--
\ref{four:maximal:sec:comparison}.  First we obtain $B$ and its outward
collars.  Next we prove a lower bound for $t$, followed by height,
boundary-flux and range estimates for every possible continuation
solution.  Local regularity and an unrestricted scalar solve then
define the compact continuation map.  Only after existence and
fixed-$N$ exhaustion do we use the end calculation of
Lemma~\ref{b:transition:end}.  The final flux estimate proves the
displayed energy bound.  This order separates the polynomial estimates
needed for the mass limit from arbitrary fixed-$N$ regularity constants.

The trace $\tr_gK$ vanishes throughout this construction.  Its
positive regularizing parameter is $\eta_N=\ell^{3/2}$; this parameter
is distinct from the trace.
The scalar weight is denoted by $L_0=e^{2t}l(t)$.
First fix the strict geometry and its decay exponents, then $\eps>0$ and
a sufficiently large $N$.  Exhaust the outer radius at fixed $N$, send
$N\to\infty$ using polynomial constants only, then send $\eps\downarrow0$
and remove the strict approximation.  Arbitrary fixed-$N$ regularity
constants do not enter the final mass limit.


\section{Strict data and the geometric inner boundary}
\label{four:maximal:sec:preparation}

This section produces the strict data and the inner boundary used in the
deformation.  The approximation is controlled in ADM energy and in the
infimum over all enclosing areas.  The geometric construction also supplies
a comparison principle for closed sets described by smooth exterior
supports; this will be used before any derivative estimates for the
deformation are available.

\subsection{Strict conformal approximation}

The spacetime Poisson construction of
Jaracz~\cite{Jaracz2025StrictDEC} is the model for this reduction.
Here the conformal exponents are four-dimensional, the inner Neumann
datum is nonzero, and we prove preservation of maximality and comparison
for the full enclosing-area infimum.

Decrease the decay exponent, if necessary, so that \(1<q<2\), and fix
\(0<\delta<q-1\).  Extend the end radius to a positive smooth function
\(r\ge1\) on the exterior.  We use the enclosing-area infimum
\(\mathcal A_g(S)\) defined above.

\begin{proposition}[Reduction to strict data]
\label{four:maximal:prop:strict-data}
The data of Theorem~\ref{four:maximal:thm:main} admit smooth approximations
\((g_j,K_j)\) on the same exterior such that
\begin{equation}
\begin{aligned}
 \tr_{g_j}K_j&=0,&
 \mu_j-|J_j|_{g_j}&\ge c_j r^{-4-\delta},\\
 \Theta_{K_j}(S)&<0,&
 R_{g_j}&=O(r^{-4-\delta}),
 \qquad c_j>0.
\end{aligned}
\label{four:maximal:prep:strict-properties}
\end{equation}
They have the prescribed \(O_2(r^{-q})\) metric decay and
\(O_1(r^{-1-q})\) tensor decay, and all their coordinate derivatives
are bounded uniformly over end patches of unit size, with constants
allowed to depend on \(j\) and the derivative order.  Moreover,
\begin{equation}
 E_{g_j}\longrightarrow E_g,
 \qquad \mathcal A_{g_j}(S)\longrightarrow\mathcal A_g(S).
\label{four:maximal:prep:approx-limits}
\end{equation}
Consequently, it suffices to prove
\(E_g\ge\tfrac12(\mathcal A_g(S)/\omega_3)^{2/3}\) for such strict maximal
data on this exterior.  Neither vacuum nor marginality is required of
the intermediate approximations.
\end{proposition}

\begin{proof}
Choose a positive smooth function \(b\ge|K|_g\) which equals a
constant multiple of \(r^{-1-q}\) sufficiently far out.  We first solve
\begin{equation}
 -\Delta_g\phi
   =b\sqrt{|\dd\phi|_g^2+r^{-6}}+r^{-4-\delta},
 \qquad \partial_\nu\phi=-1\text{ on }S,
 \qquad \phi\longrightarrow0\text{ at infinity},
\label{four:maximal:prep:phi-equation}
\end{equation}
where \(\nu\) points into the exterior.  In particular, the normal
derivative in this equation is the negative of the domain's outward
normal derivative at its inner boundary.

Here are the existence and decay details.  On the compact truncation
\(M_R\), impose \(\phi=0\) on the outer coordinate sphere and replace
\(b\) on the right by \(\lambda b\), \(0\le\lambda\le1\).
The linearization, with homogeneous boundary data, is
\[
 v\longmapsto-\Delta_gv
 -\lambda b\frac{\langle\dd\phi,\dd v\rangle_g}
                        {\sqrt{|\dd\phi|_g^2+r^{-6}}}.
\]
Its drift is bounded by \(b\).  Its kernel is zero by the strong
maximum principle and the boundary point lemma, since the outer
Dirichlet face is nonempty and the inner face has homogeneous normal
derivative.  The mixed operator has Fredholm index zero, by deformation
to the mixed Laplacian, and hence is invertible.  The two boundary
components are disjoint, smooth, and compact, so ordinary scalar
elliptic estimates apply without a junction between boundary types;
see \cite[Chapters~3, 6, and~9]{GilbargTrudinger2001} and
\cite[Theorems~7.3 and~15.2--15.3]{AgmonDouglisNirenberg1959}.

For fixed \(R\), the global mixed \(W^{2,p}\) estimate, \(p>4\),
and interpolation give
\begin{equation}
 \|\phi\|_{W^{2,p}(M_R)}
 \le C_{R,p}\bigl(1+\|\dd\phi\|_{L^p(M_R)}
                         +\|\phi\|_{L^p(M_R)}\bigr)
 \le C'_{R,p}(1+\|\phi\|_\infty).
\label{four:maximal:prep:phi-estimate}
\end{equation}
These constants are uniform in \(\lambda\).  If the suprema were
unbounded, divide by their sup norms and take a subsequence.  Compact
embedding of \(W^{2,p}\) in \(C^{1,\alpha}\), with a smaller
\(\alpha>0\), gives a nonzero limit \(v\) satisfying
\[
 -\Delta_gv=\lambda_*b|\dd v|_g,
 \qquad \partial_\nu v=0\text{ on }S,
 \qquad v=0\text{ on the outer face}.
\]
This is a homogeneous linear equation with a bounded measurable drift:
set the drift to \(\lambda_*b\grad v/|\dd v|\) where
\(\dd v\ne0\), and to zero otherwise.  The maximum principle and
boundary point lemma again exclude a nonzero solution.  Thus
\eqref{four:maximal:prep:phi-estimate} is bounded, and Schauder estimates and
bootstrapping give closedness in the continuity argument.  The
\(\lambda=0\) mixed Poisson problem is solvable, so the implicit
function theorem and compactness give a smooth solution for
\(\lambda=1\).  It is nonnegative: a negative minimum is excluded
in the interior by \(-\Delta\phi>0\), on the outer face by its
zero value, and on the inner face by \(\partial_\nu\phi=-1\).

We next bound these solutions independently of \(R\).  On the end
consider
\[
 T(r)=r^{-2}(1-Ar^{-\delta}).
\]
Choose a fixed \(A>0\), then a large fixed \(r_0\), so that \(T>0\)
and
\begin{equation}
 (-\Delta_g-b|\dd\,\cdot\,|)T
 \ge c\,r^{-4-\delta}\qquad(r\ge r_0).
\label{four:maximal:prep:radial-barrier}
\end{equation}
Indeed,
\(-\Delta_\R T=A\delta(2+\delta)r^{-4-\delta}\), while both
the metric error and \(b|\dd T|\) are \(O(r^{-4-q})\).
Also
\[
 -\Delta_g\phi-b|\dd\phi|
 \le b r^{-3}+r^{-4-\delta}\le C r^{-4-\delta}.
\]
Comparison on \(\{r_0\le r\le R\}\), including the zero outer
value, therefore yields
\begin{equation}
 0\le\phi_R(x)\le C(1+m_R)r(x)^{-2},
 \qquad m_R=\max_{\{r\le r_0\}}\phi_R.
\label{four:maximal:prep:uniform-tail}
\end{equation}
If \(m_R\) were unbounded along expanding truncations, divide by
\(m_R\).  Local versions of \eqref{four:maximal:prep:phi-estimate}, including
the inner boundary estimates, give a nonnegative limit attaining
value one in the fixed core.  It solves
\(-\Delta v=b|\dd v|\), has homogeneous inner normal derivative,
and is \(O(r^{-2})\) by \eqref{four:maximal:prep:uniform-tail}.  Its positive
maximum is therefore attained at a finite point.  The strong maximum
principle or the boundary point lemma excludes it.  We have a uniform
core bound, and diagonal compactness gives a smooth solution of
\eqref{four:maximal:prep:phi-equation}.

Rescale annuli by their radii.  The bound \(\phi=O(r^{-2})\),
the scaled metric bounds furnished by \(g-\delta=O_2(r^{-q})\),
and the linear gradient growth in \eqref{four:maximal:prep:phi-equation} first
give scaled \(W^{2,p}\) estimates by interpolation.  They give
H\"older control of the gradient; the right side then has the
scaled H\"older control needed for Schauder estimates.  Hence
\begin{equation}
 \phi=O_2(r^{-2}).
\label{four:maximal:prep:phi-decay}
\end{equation}
The source in \eqref{four:maximal:prep:phi-equation} is
\(O(r^{-4-q})+O(r^{-4-\delta})\), and is integrable.  The
divergence theorem shows that
\[
 I_\phi=\lim_{R\to\infty}\int_{S_R}\partial_{\nu_g}\phi\dd A_g
\]
exists.  Replacing this by its Euclidean-coordinate flux changes it
by \(O(R^{-q})\), which tends to zero.

For small \(s>0\) set
\(g_s=e^{2s\phi}g\) and \(K_s=e^{s\phi}K\).
The trace remains zero.  The connection change is
\[
 \Gamma(g_s)^k_{ij}-\Gamma(g)^k_{ij}
 =s(\delta_i^k\phi_j+\delta_j^k\phi_i-g_{ij}\phi^k).
\]
Contracting it against the scaled trace-free tensor, and using the
four-dimensional scalar-curvature formula, gives
\begin{equation}
\begin{split}
 e^{2s\phi}\mu_s
   &=\mu-3s\Delta_g\phi-3s^2|\dd\phi|^2,\\
 e^{s\phi}J_s
   &=J+3sK(\grad\phi,\cdot),\\
 \Theta_{K_s}(S)
   &=e^{-s\phi}\bigl(\Theta_K(S)+3s\partial_\nu\phi\bigr).
\end{split}
\label{four:maximal:prep:conformal-constraints}
\end{equation}
The last coefficient is three because the boundary has dimension
three.  A covector norm contributes one more factor \(e^{-s\phi}\),
so the first two identities and \(\mu=J=0\) imply
\begin{align}
 e^{2s\phi}(\mu_s-|J_s|_{g_s})
 &\ge 3s\bigl[-\Delta\phi-|K||\dd\phi|-s|\dd\phi|^2\bigr]
 \notag\\
 &\ge3s\bigl[r^{-4-\delta}-s|\dd\phi|^2\bigr].
\label{four:maximal:prep:strict-gap}
\end{align}
The ratio \(|\dd\phi|^2/r^{-4-\delta}\) is bounded, by
\eqref{four:maximal:prep:phi-decay} and \(\delta<1\).  Thus small \(s>0\)
gives a positive gap of the asserted order.  The boundary expansion
is \(-3s e^{-s\phi}<0\).  Moreover,
\[
 R_{g_s}=e^{-2s\phi}
 \bigl(|K|^2-6s\Delta\phi-6s^2|\dd\phi|^2\bigr)
       =O(r^{-4-\delta}),
\]
because \(2+2q>4+\delta\).  The original metric and tensor decay
is preserved, since \(q<2\).

Only \(-6s\partial_i\phi\) survives from the conformal change
in the ADM integrand.  Products with a metric error or with \(\phi\)
have vanishing integrated flux by \eqref{four:maximal:prep:phi-decay}.  Thus
\begin{equation}
 E_{g_s}=E_g-\frac{s}{\omega_3}I_\phi\longrightarrow E_g.
\label{four:maximal:prep:conformal-mass}
\end{equation}
Since \(\phi\) is bounded, the comparison
\(e^{-2s\|\phi\|_\infty}g\le g_s\le
e^{2s\|\phi\|_\infty}g\) holds on the entire exterior.  Its
three-dimensional area comparison holds for every enclosing cut,
and therefore for their infima.  In particular,
\(\mathcal A_{g_s}(S)\to\mathcal A_g(S)\), with no compactness assumption on
minimizing cuts.

\paragraph{Smoothing the end.}
The conformal step has produced strict dominance and convergence of energy
and enclosing-area infima.  The differentiated end estimates will also
need bounds of every fixed derivative order on patches of unit size.
The original \(O_2/O_1\) hypotheses do not supply those bounds; we obtain
them by smoothing each fixed strict pair sufficiently far out.

Fix one of these strict pairs and temporarily denote it by \((g,K)\).
Let \(h=g-\delta\) in the end chart.  Convolve \(h\) and \(K\)
with a nonnegative smooth kernel of integral one supported in a unit
ball.  Interpolate to the original data with a cutoff which is zero
for \(r\le R_{\rm sm}\), one for \(r\ge2R_{\rm sm}\), and has
derivatives of order \(i\) bounded by \(C_iR_{\rm sm}^{-i}\).
Call the resulting metric and tensor \(\widetilde g,\widetilde K\),
and project the tensor to its \(\widetilde g\)-trace-free part:
\[
 K^{\rm sm}=\widetilde K-\tfrac14
                (\tr_{\widetilde g}\widetilde K)\widetilde g.
\]
For sufficiently large \(R_{\rm sm}\), the metric remains positive
definite.  The available original derivatives give
\[
 h*\kappa-h=O(r^{-q-1}),\qquad
 D(h*\kappa-h)=O(r^{-q-2}),\qquad
 K*\kappa-K=O(r^{-q-2}).
\]
Here \(D\) denotes coordinate differentiation only in this paragraph.
For the trace projection, use the original trace-free identity in the
form
\[
 \tr_\delta K=(\delta^{ij}-g^{ij})K_{ij}
                       =O_1(r^{-1-2q}).
\]
Convolution and interpolation preserve this first-derivative bound.
Moreover
\[
 \tr_{\widetilde g}\widetilde K
 =\tr_\delta\widetilde K
      +(\widetilde g^{ij}-\delta^{ij})\widetilde K_{ij}
 =O_1(r^{-1-2q}).
\]
Differentiating the last product uses only the available first derivatives
of the metric and tensor.  Thus the projection changes the tensor by
\(O_1(r^{-1-2q})\), without requiring a second derivative of the original
\(K\).

For trace-free data with these decay bounds, the constraint pair has
the flat linear principal expressions
\begin{equation}
 (\mu,J)
 =\left(\tfrac12(\partial_i\partial_jh_{ij}
                      -\Delta_\R h_{ii}),\,
                   \partial_jK_{ij}\right)
     +O(r^{-2-2q}).
\label{four:maximal:prep:linear-constraints}
\end{equation}
The scalar error consists of \(hD^2h\), \((Dh)^2\), and
\(|K|^2\); the momentum error consists of \(hDK\) and \((Dh)K\).
The displayed linear operators commute with convolution.  Cutoff
derivatives multiply the preceding convolution differences and
contribute \(O(r^{-q-3})\).  The trace projection changes the
momentum by \(O(r^{-2-2q})\).  Hence the new constraint pair is the
interpolated average of the old pair, up to
\begin{equation}
 O(r^{-2-2q}+r^{-q-3})=o(r^{-4-\delta}).
\label{four:maximal:prep:smoothing-errors}
\end{equation}
Replacing the old momentum norm by its Euclidean norm, and the averaged
Euclidean momentum norm by the new metric norm, each costs
\(O(r^{-2-2q})\): both metrics differ from the Euclidean metric by
\(O(r^{-q})\), while both momenta are \(O(r^{-2-q})\).
The triangle inequality and positivity
of the kernel therefore preserve a fixed positive fraction of the strict
dominance gap once \(R_{\rm sm}\) is sufficiently large.  This radius is
chosen separately for each fixed strict pair, so the little-o errors are
smaller than its particular positive gap coefficient.  The two needed
exponent inequalities are
\[
 2q-2-\delta>0,\qquad q-1-\delta>0.
\]
The same calculation shows that the scalar curvature is still
\(O(r^{-4-\delta})\).

Beyond \(2R_{\rm sm}\), derivatives of the convolution kernel
bound every higher derivative on unit patches, using the original
zeroth, first, and second derivative bounds.  The transition region
is smooth and compact.  At infinity the metric's first derivative
changes by \(O(r^{-q-2})\), whose integrated ADM contribution is
\(O(r^{1-q})\to0\).  Thus this smoothing leaves the energy exactly
unchanged.  As \(R_{\rm sm}\to\infty\), its uniform metric ratio
tends to one, so the enclosing infimum converges as well.  Choosing
first \(s\downarrow0\), and then a sufficiently large smoothing
radius for each \(s\), proves \eqref{four:maximal:prep:strict-properties} and
\eqref{four:maximal:prep:approx-limits}.
\end{proof}

For the remainder of the proof fix strict data from
Proposition~\ref{four:maximal:prop:strict-data}, and write
\(A_*=\mathcal A_g(S)\).  Constants may depend on this fixed approximation.
Write \(c_0>0\) for a fixed constant such that
\(\mu-|J|\ge c_0r^{-4-\delta}\) for these data.

\subsection{Threshold regions and outward collars}

We use the following part of the total trapped-region theorem:
on a complete smooth asymptotically flat initial data set of spatial
dimension \(2\) through \(7\), the nonempty union of bounded regions
with smooth boundary and outward expansion at most zero has a smooth
embedded outermost stable MOTS as its outer frontier.  We take its
closed, filled representative.  Filling a bounded complementary
component deletes boundary components and preserves admissibility.
This is the existence and regularity conclusion of
\cite[Section~4.3, Theorem~4.6]{AnderssonEichmairMetzger2011}; it requires no energy
condition.  We do not use the additional topology conclusions of that
theorem.

\begin{lemma}[Threshold selection and collars]
\label{four:maximal:lem:threshold-collars}
There is a number \(c\) with \(\max_S\Theta_K<c<0\), and a
closed filled compact region \(\mathcal D_c\), containing a smooth
compact filling of \(S\) and an exterior collar of \(S\), with the
following properties.
Its boundary \(B\) is a smooth embedded stable threshold surface,
\(\Theta_K(B)=c\).  It belongs to an increasing family
\(\mathcal D_b\), \(\max_S\Theta_K<b<0\), of total regions
at threshold \(b\), all lying in a fixed compact set, and
\begin{equation}
 \mathcal D_{b}\longrightarrow\mathcal D_c
       \quad\text{in Hausdorff distance as }b\downarrow c.
\label{four:maximal:prep:right-continuity}
\end{equation}
The exterior \(\Omega\) of \(\mathcal D_c\) is connected and
one-ended.  Every smooth cut enclosing \(B\) has \(g\)-area at
least \(A_*\).  Each component of \(B\) has a smooth outward
collar with leaf coordinate \(s\ge0\), \(s=0\) at \(B\),
\(|\dd s|>0\), whose leaves satisfy
\begin{equation}
 \Theta_K(\{s=\text{constant}\})\ge c,
 \qquad \nu_s=\frac{\grad s}{|\dd s|}.
\label{four:maximal:prep:outward-collars}
\end{equation}
\end{lemma}

\begin{proof}
The given diffeomorphism of the exterior permits a smooth compact
filling behind \(S\).  Extend \(g,K\) smoothly over it; the
extended metric can be chosen positive definite by interpolation
with any interior metric away from a collar.  This extension is used
only in the present geometric construction.

Fix a large radius \(R_0\) beyond which every coordinate sphere has
\(\Theta_K=3/r+O(r^{-1-q})>0\).  Outside a still larger fixed
radius, smoothly cut the metric to the Euclidean metric and the
unshifted tensor to zero.  Dilated cutoffs preserve the bounds
\(g-\delta=O_2(r^{-q})\), \(K=O_1(r^{-1-q})\), and consequently
the positive expansion of all spheres in the cutoff annuli.  Denote
these complete auxiliary data by \((g^{\rm aux},K^{\rm aux})\).
Choose a fixed nonnegative smooth cutoff \(\xi\), equal to one
on the filling and throughout a large neighborhood of \(\{r\le
R_0\}\), and zero far out.  For
\(\max_S\Theta_K<b<0\), use the tensor
\[
 K_b^{\rm aux}=K^{\rm aux}-\tfrac b3\xi g^{\rm aux}.
\]
Where \(\xi=1\) and the data are uncut, the exact identity is
\begin{equation}
 \Theta_{K_b^{\rm aux}}(\Sigma)=\Theta_K(\Sigma)-b.
\label{four:maximal:prep:threshold-shift}
\end{equation}
In the cutoff region the added tangential trace is \(-b\xi\ge0\),
so all the distant spheres still have positive expansion, uniformly
in \(b\).

A smooth weakly trapped bounded region cannot meet those distant
spheres: at a point where its boundary reaches its greatest radius,
its tangent plane and outward normal agree with those of the
coordinate sphere, and its mean curvature is at least the sphere's
mean curvature.  The tensor traces there agree.  This contradicts
nonpositive expansion.  The same argument applies to the outer
frontier furnished by the total-region theorem.  Thus all total
regions are contained in the uncut part of the construction.

For every indicated \(b\), strictness at \(S\) gives a short
outward collar whose leaves have expansion less than \(b\).
The filling together with this collar is a nonempty trapped region
for \(K_b^{\rm aux}\).  The theorem therefore gives a total closed
region \(\mathcal D_b\) with smooth stable frontier.  Its frontier
lies strictly outside \(S\), in the original data, and has
\(\Theta_K=b\) by \eqref{four:maximal:prep:threshold-shift}.
We used one filling and one set of cutoffs for the entire interval
of \(b\).  Since \(\xi\ge0\), increasing \(b\) decreases
the shifted expansion on every candidate hypersurface.  It follows
that
\begin{equation}
 b_1<b_2\quad\Longrightarrow\quad
                 \mathcal D_{b_1}\subset\mathcal D_{b_2}.
\label{four:maximal:prep:threshold-monotonicity}
\end{equation}

Choose a compact enlargement \(Q\) containing all these regions.
Lemma~\ref{n3:domains:right-continuity} applies to the increasing family
of closed subsets \(E_b=\mathcal D_b\) of this fixed compact metric
space, with \(\max_S\Theta_K<b<0\).  The preceding construction proves
both the common compact containment and monotonicity; every member is
nonempty.  Choose \(c\) outside the lemma's countable exceptional set.
Its conclusion is exactly \eqref{four:maximal:prep:right-continuity}.
Filling bounded complementary
pockets ensures that the remaining exterior is the connected
component containing the sole end.  Any cut enclosing its entire
boundary \(B\) also encloses \(S\), proving the area assertion.

It remains to construct the collars.  Apply
Lemma~\ref{n3:domains:stable-collar} with ambient dimension four,
\(Q=K\), threshold \(c\), and each connected component of \(B\).
The total-region theorem gives a compact smooth stable frontier;
its principal eigenvalue is nonnegative
\cite[Proposition~5.1]{AnderssonMarsSimon2008}.  On an open neighborhood
of this frontier the auxiliary cutoff \(\xi\) equals one and the data
are uncut, so the relevant shifted tensor is exactly \(K-(c/3)g\).
Its expansion differs from \(\Theta_K\) by the constant \(c\), with
zero derivative under graph variations.  Replacing one component by
a sufficiently short outward graph of expansion at most \(c\)
would strictly enlarge the total region while preserving every other
component.  Thus the local maximality hypothesis of the collar lemma
holds in this application.

For clarity, its zero-mode construction uses the expansion
\(\mathcal E(v)\) of the outward graph on this one component and
\(L=D\mathcal E(0)\).  The common augmented inverse specializes to
\begin{equation}
 (v,a)\longmapsto\left(Lv-a,\ \int_Bv\dd A_g\right),
 \qquad C^{2,\alpha}(B)\times\R
             \longrightarrow C^{0,\alpha}(B)\times\R.
\label{four:maximal:prep:augmented-inverse}
\end{equation}
Here the integrals are over the chosen component.  The positive
principal and adjoint eigenfunctions, their simple kernels, and the
positive pairing used to invert this map are those in
\cite[Section~4, Lemma~4.1]{AnderssonMarsSimon2008}, as in the common
proof.  That proof gives strictly larger expansion on every positive
outward leaf, including the zero-eigenvalue case; stability alone
would not give this conclusion.  Finitely many components permit
disjoint shortened collars, proving
\eqref{four:maximal:prep:outward-collars}.
\end{proof}

\subsection{Exterior supports of nonsmooth sets}

The threshold regions control smooth surfaces.  The graph-height limit
will give only a closed sublevel set, so we need the same comparison for
smooth surfaces that touch such a set from outside.  The next argument
transports the supporting tangent plane exactly; this prevents an
uncontrolled curvature of the support from entering the error.

For a closed set \(A\), a smooth exterior support at
\(x\in\partial A\) is a regular level \(\{\psi=0\}\) in a
neighborhood of \(x\), where \(\psi(x)=0\), \(\dd\psi(x)\ne0\),
and \(A\subset\{\psi\le0\}\) locally.  Its normal is
\(\nu_\psi=\grad\psi/|\dd\psi|\), and its expansion is
\begin{equation}
 \Theta_K[\psi]
 =\frac{\tr_g\Hess\psi-\Hess\psi(\nu_\psi,\nu_\psi)}
              {|\dd\psi|}
      +\tr_gK-K(\nu_\psi,\nu_\psi).
\label{four:maximal:prep:support-expansion}
\end{equation}
Only supports that exist are tested; no regularity is assumed of
\(A\) itself.

\begin{lemma}[Enclosure from exterior supports]
\label{four:maximal:lem:support-enclosure}
Let \(A\) be a compact closed set containing \(\mathcal D_c\).
Suppose that every smooth exterior support to its frontier has
expansion at most \(k'\), where \(c\le k'<0\).  Then
\[
 A\subset\mathcal D_{c'}\qquad\text{for every }k'<c'<0.
\]
All expansions at the frontier of \(A\) are computed in the
original strict exterior data.
\end{lemma}

\begin{proof}
First, every set under consideration lies inside one fixed large
sphere.  Otherwise the coordinate sphere at its maximum radius is
an exterior support with positive expansion, contrary to \(k'<0\).
Choose a compact enlargement containing this sphere and a small
metric neighborhood of it.  All constants below depend only on
the smooth geometry and tensor on this enlargement.  Distances
may be computed in the complete smooth filling, since the contact
points and sufficiently short segments used below lie in the
original exterior: \(A\) contains the filling and a collar of
\(S\).

For \(z>0\), put \(A_z=\{y:\dist(y,A)\le z\}\).
Apply Lemma~\ref{n3:domains:parallel-support} in ambient dimension
four, with its set \(D=A\), tensor \(Q=K\), and threshold
\(k'\).  The compact neighborhood and short segments required by
that local lemma are precisely those just fixed.  The definition
\eqref{four:maximal:prep:support-expansion} is its normalized
expansion.  It follows that every exterior support to \(A_z\)
has expansion at most
\begin{equation}
 k'+Cz,
\label{four:maximal:prep:parallel-support-bound}
\end{equation}
where \(C\) is independent of that support's curvature.

The exact geometric feature of this application is worth recording.
For a nearest pair \(x\in A\), \(x'\in\partial A_z\), let
\(P:T_xM\to T_{x'}M\) be parallel transport along the short
minimizing segment.  The Jacobi construction in the common proof gives
\begin{equation}
 F_z(x)=x',\qquad \dd F_z|_x=P,
 \qquad |\nabla\dd F_z|_x\le Cz.
\label{four:maximal:prep:exact-transport}
\end{equation}
If \(\psi\) defines the support at \(x'\), then
\(\psi\circ F_z\) supports \(A\) at \(x\), and
\begin{equation}
\begin{split}
 \Hess(\psi\circ F_z)_x(X,Y)
  ={}&\Hess\psi_{x'}(PX,PY)\\
    &+\dd\psi_{x'}\bigl((\nabla\dd F_z)_x(X,Y)\bigr).
\end{split}
\label{four:maximal:prep:hessian-pullback}
\end{equation}
The exact tangent isometry preserves the normalized Hessian trace;
only the \(O(z)\) second-map derivative and the \(C^1\) change
of \(K\) enter the expansion error.  Thus no term multiplying
\(\Hess\psi\) has been discarded.  This local transport step
uses compact ambient bounds, while the following enclosure argument
still uses the auxiliary total regions specific to the maximal end.

Fix \(k'<c'<0\), and take \(\beta>0\) below the preceding
injectivity and collar scales, so small that
\(k'+C\beta<c'\).  Smooth approximation of the distance function,
followed by a regular-value choice, gives a smooth compact region
\(U\) with
\[
 A_{\beta/4}\subset\operatorname{int}U,
 \qquad U\subset A_{3\beta/4}.
\]
For example, approximate the distance uniformly within
\(\beta/32\), and choose a regular value between
\(3\beta/8\) and \(5\beta/8\).  Choose a smooth nonnegative
function \(\zeta\) equal to one near \(\partial U\) with
\(\supp\zeta\subset\{\dist(\cdot,A)<\beta\}\).
For sufficiently large finite \(\Lambda\), replacing the
auxiliary threshold tensor by
\[
 K_{c'}^{\rm aux}-\Lambda\zeta g^{\rm aux}
\]
makes \(\partial U\) strictly trapped: the added term decreases
its expansion by \(3\Lambda\).  Apply the total-region theorem
to these smoothly modified complete data.  Let \(\widetilde{\mathcal
D}\) be the resulting closed filled total region and
\(\widetilde B\) its smooth frontier.  It contains \(U\).
All cutoffs were placed outside the fixed enlargement above;
outside the support of \(\zeta\), large spheres retain positive
expansion.  The maximum-radius argument confines
\(\widetilde B\) to the uncut exterior.  There it satisfies
\begin{equation}
 \Theta_K(\widetilde B)=c'+3\Lambda\zeta\ge c'.
\label{four:maximal:prep:modified-frontier}
\end{equation}

Let \(z=\dist(A,\widetilde B)\), the global minimum distance.
Because \(\widetilde{\mathcal D}\) contains
\(A_{\beta/4}\), one has \(z\ge\beta/4>0\).
If \(z<\beta\), every path of length less than \(z\) starting
in \(A\) remains on the inner side of \(\widetilde B\): its
first crossing would contradict the definition of \(z\).
Taking closures proves \(A_z\subset\widetilde{\mathcal D}\).
At a pair attaining the distance, \(\widetilde B\) is therefore
an exterior support to \(A_z\).  Equations
\eqref{four:maximal:prep:parallel-support-bound} and
\eqref{four:maximal:prep:modified-frontier} imply
\[
 c'\le\Theta_K(\widetilde B)\le k'+Cz
                  \le k'+C\beta<c',
\]
a contradiction.  Thus \(z\ge\beta\).  Since this is the
global minimum, \(\zeta\) vanishes at every point of
\(\widetilde B\), including when the minimum equals \(\beta\).
It follows that \(\Theta_K(\widetilde B)=c'\) everywhere.
Its filled inner region is an admissible threshold-\(c'\) region
for the original auxiliary data, so
\(A\subset\widetilde{\mathcal D}\subset\mathcal D_{c'}\),
as required.
\end{proof}

\subsection{Fixed collar data}

Set \(b_*=c\).  On the disjoint collars supplied by
Lemma~\ref{four:maximal:lem:threshold-collars}, choose a nonpositive smooth
compactly supported function \(C\), equal to \(-\kappa s\)
near each component of \(B\), with \(\kappa>0\) sufficiently
small.  Indeed, multiply \(-\kappa s\) by nonnegative collar
cutoffs.  The strict gap has a positive minimum on their compact
support, so \(\kappa\) can be chosen small enough that, for a
fixed smooth positive function \(\rho\) equal to a positive
multiple of \(r^{-4-\delta}\) on the end,
\begin{equation}
 |\dd C|+\rho\le\mu-|J|.
\label{four:maximal:prep:offset-slack}
\end{equation}
For instance first choose \(\rho\le(c_0/4)r^{-4-\delta}\),
then bound \(|\dd C|\) by half the remaining strict gap on its
support.  These data, the collars, the threshold \(b_*\), and
the region \(\Omega\) are fixed before choosing any of the
deformation parameters below.

\section{The scalar-curvature deformation}\label{four:maximal:sec:deformation}

We now fix the strict data and the exterior $\Omega$ constructed in
Section~\ref{four:maximal:sec:preparation}.  Thus $B=\partial\Omega$ has expansion
$H_B+\tr_BK=b_*<0$, the tensor is maximal, and the fixed functions $C\le0$
and $\rho>0$ satisfy
\begin{equation}\label{four:maximal:def:slack}
 |dC|+\rho\le\mu-|J|.
\end{equation}
The normal $\nu$ at $B$ points into $\Omega$.  The compactly supported
function $C$ vanishes on $B$ and equals $-\kappa s$ in its fixed collars.
Throughout this section the base dimension is four and a transverse block
has dimension three.

\subsection{Variables and the exact scalar identity}

The pointwise identities used here are those of
Section~\ref{four:sec:system}.  We retain their variables and
four-dimensional density conventions below.  Their proofs are local
and impose no support condition on $K$; maximality will enter through
$\tr_gK=0$ in the equations and the estimates for the tensor tail.

Fix $0<\eps<1$.  The integer $N\ge4$ is an auxiliary parameter.  Fix once
and for all a smooth nonincreasing function $\vartheta:\R\to[0,1]$ equal
to one on $(-\infty,0]$ and zero on $[1,\infty)$, and set
\begin{equation}\label{four:maximal:def:warp}
 \begin{gathered}
 p(t)=N\vartheta(Nt),\qquad
 l(t)=\exp\left(\int_0^t p(s)\dd s\right),\qquad
 L_0(t)=e^{2t}l(t),\\
 \ell=e^{-\eps N},\qquad \eta_N=\ell^{3/2}.
 \end{gathered}
\end{equation}
In particular $0\le p\le N$, $|p'|\le\|\vartheta'\|_\infty N^2$,
$l(t)=e^{Nt}$ for $t\le0$,
and $l\le e$ everywhere.  A symbol $\Pi_N$ denotes a positive polynomial
bound in $N$; its coefficients may depend on the fixed strict data,
collars, cutoffs, $\eps$, and subsequently fixed small constants, but not
on a solution, the outer radius, or a homotopy parameter.  Its value may
increase from one estimate to the next.  Arbitrary constants at fixed
$N$ will instead be denoted by $C(N)$.

On $\Omega_R=\Omega\cap\{r<R\}$ the unknowns are $f,Z$.  Define $t$
implicitly and the remaining variables explicitly by
\begin{align}
 \sigma&=|\grad f|,& D&=1+l(t)^2\sigma^2,&d&=D^{-1},
 & Z&=t+\tfrac16\log D,\label{four:maximal:def:implicit}\\
 w&=\frac{l\grad f}{\sqrt D},&a&=|w|,&v&=a^2=1-d,
 &\chi&=\frac{3d}{3+pv},\label{four:maximal:def:axial-variables}\\
 h&=\eta_N f,&U&=l\sqrt D,&u&=L_0\sqrt D.
 \label{four:maximal:def:weights}
\end{align}
At fixed $\sigma$, the derivative of the right side of
\eqref{four:maximal:def:implicit} with respect to $t$ is $1+pv/3>0$.  As $t\to-\infty$
it tends to $-\infty$, and as $t\to\infty$ it tends to $\infty$.
Consequently $t=t(Z,df)$ exists uniquely and smoothly for all finite
arguments.  Notice that $t\ge-\eps$ implies $\ell\le l\le e$.

For $\sigma>0$ put $e_f=\grad f/\sigma$ and
$A_j=\Id+(j-1)e_f\otimes e_f$.  The two matrices used below have the
globally smooth formulas
\begin{equation}\label{four:maximal:def:matrices}
 \Ad=\Id-w\otimes w,\qquad
 \Achi=\Id-\frac{3+p}{3+pv}w\otimes w,
 \qquad \frac{3d}{3+N}\le\chi\le d\le1.
\end{equation}
They are positive definite; the displayed formulas remove the apparent
ambiguity of the axis when $df=0$.  We identify vectors and covectors
using $g$, write $|\xi|_A^2=A(\xi,\xi)$, and put
\begin{equation}\label{four:maximal:def:graph-fields}
 \begin{gathered}
 \gbar=g+l^2df\otimes df,\quad
 \ghat=e^{2t}\gbar,\quad
 H^f=\frac{l}{\sqrt D}\Hess f,\quad S_f=K+H^f,\\
 V=u\Achi\bigl(3\grad Z+K(w,\cdot)\bigr).
 \end{gathered}
\end{equation}
The trace equation will always be
\begin{equation}\label{four:maximal:def:trace}
 \tr_{\Achi}S_f=F.
\end{equation}
Here $F$ is the prescribed right side, including its dependence on $f$,
$Z$, and $df$ when a system is specified.

For symmetric tensors recall
$P_A=A-(\tr A)g$ and
$\mathsf{B}(A_1,A_2)=A_1:A_2-(\tr A_1)(\tr A_2)$.
The constraint quantities are
$\mu=(R_g-\mathsf{B}(K,K))/2$ and $J=\divg P_K$.
At a point with $\sigma>0$, decompose $S_f$ relative to
$e_f^\perp\oplus\R e_f$ into its transverse symmetric block $T$, mixed
covector $M$, and axial entry $j$.  The letter $M$ in the formulas below
denotes only this block.  Put
\begin{equation}\label{four:maximal:def:blocks}
 x=\partial_{e_f}t,\qquad y=(\grad t)_\perp,
 \qquad T^0=T-\tfrac13(F-\chi j)\Id_\perp.
\end{equation}
The trace equation gives $\tr T=F-\chi j$.

\begin{proposition}[Scalar identity]\label{four:maximal:prop:scalar-identity}
For smooth $f,t$ satisfying \eqref{four:maximal:def:trace},
\begin{equation}\label{four:maximal:def:scalar-identity}
 \frac12 e^{2t}R_{\ghat}
 =\mu+J(w)+\cT+w(F)-F\tr K-u^{-1}\divg V,
\end{equation}
where
\begin{align}\label{four:maximal:def:quadratic-form}
 \cT={}&\tfrac12|T^0|^2+|M+(p+1)ay|^2-v|y|^2
       +3(p+1)(|y|^2+dx^2)\notag\\
 &-\tfrac13(F-\chi j)^2+\chi j^2-\chi Fj+F^2
       +2(p+1)\chi jax.
\end{align}
This expression extends smoothly across $df=0$.
\end{proposition}

\begin{proof}
This is Proposition~\ref{four:sys:scalar-proposition}, with its axis
$e=e_f$ and the same graph variables and density normalization.
Its arbitrary trace $\tau$ is zero here.

We record the differential relations needed in the boundary and drift
calculations.  Differentiating the implicit definition of $Z$ gives
\begin{equation}\label{four:maximal:def:Z-differential}
 3\,dZ=(3+pv)dt+H^f(w,\cdot).
\end{equation}
If $k=-H^f-p(dt\otimes w^\flat+w^\flat\otimes dt)$ and $Q=K-k$, then
\begin{align}
 V/u&=P_Qw+3\Ad\grad t+Fw,
                    \label{four:maximal:def:flux-rewrite}\\
 u^{-1}\divg(uw)
 &=F+(1-\chi)j-\tr_gK+2(p+1)w(t).
                    \label{four:maximal:def:w-divergence}
\end{align}
These are the identities in the proof of the cited proposition.
They follow by inserting $\tr T=F-\chi j$ and
$H^f(w,\cdot)=v\,d\log\sigma$ into its invariant formula; their
smooth tensor expressions also hold at $df=0$.  The quadratic expression
there is precisely the displayed $\cT$.
\end{proof}

\begin{lemma}[Polynomial coercivity]\label{four:maximal:lem:coercivity}
The quadratic form in \eqref{four:maximal:def:quadratic-form} is nonnegative, and
\begin{equation}\label{four:maximal:def:coercivity}
 |T|^2+|M|^2+dj^2+F^2+|dt|_{\Ad}^2\le\Pi_N\cT.
\end{equation}
At $dt=0$ it has the stronger equivalence
\begin{equation}\label{four:maximal:def:coercivity-stationary}
 \cT\asymp |T^0|^2+|M|^2+\chi j^2+F^2,
\end{equation}
with both constants absolute, independent of $N,d,\chi$.
\end{lemma}

\begin{proof}
Lemma~\ref{four:sys:coercivity-lemma} applies with the same dictionary.
For clarity, its exact square completion is
\begin{align}\label{four:maximal:def:square-completion}
 \cT={}&\tfrac12|T^0|^2+|M+(p+1)ay|^2
       +\bigl(3(p+1)-v\bigr)|y|^2\notag\\
 &+3(p+1)d\left(x+\frac{\chi aj}{3d}\right)^2
       +\tfrac23\left(F-\frac{\chi j}{4}\right)^2
       +\frac{d(48-3d)}{8(3+pv)^2}j^2.
\end{align}
The last coefficient is at least $45d/[8(3+N)^2]$ and the coefficient
of $|y|^2$ is at least two.  Undoing the shifted squares therefore
costs only a polynomial in $N$.  When $dt=0$, the normal terms instead
reduce to
\[
 \tfrac23(F-\chi j/4)^2+\chi(1-3\chi/8)j^2.
\]
The bounds $5/8\le1-3\chi/8\le1$ and $\chi^2\le\chi$ give the
absolute equivalence claimed in \eqref{four:maximal:def:coercivity-stationary}.
\end{proof}

\begin{lemma}[Boundary identity]\label{four:maximal:lem:boundary-identity}
On a hypersurface where $f$ is constant, let $H$ denote its $g$-mean
curvature for a chosen normal $\nu$, and put $P_B=\tr_BK$.  Then
\begin{align}
 V_\nu/u&=d(H+3\partial_\nu t)-H+w_\nu(F-P_B),
 \label{four:maximal:def:boundary-flux}\\
 \widehat H&=e^{-t}\sqrt d\,(H+3\partial_\nu t).
 \label{four:maximal:def:boundary-curvature}
\end{align}
The identities allow either sign of $\partial_\nu f$, including zero.
\end{lemma}

\begin{proof}
Apply Lemma~\ref{four:sys:boundary-lemma} with the chosen normal $\nu$
and trace relation \eqref{four:maximal:def:trace}.  Its formulas use
the signed component $w_\nu$ and extend smoothly to $df=0$, so they
also apply when the normal graph slope changes sign.
\end{proof}

\subsection{The equations and their homotopy}\label{four:maximal:def:homotopy}

Choose
\begin{equation}\label{four:maximal:def:end-weights}
 1<\gamma<\min(2,q-\delta/2),\qquad
 \rho_0=r^{-4-\delta},\qquad \rho_1=r^{-2\gamma}.
\end{equation}
Here $r$ is the fixed positive smooth extension of the end radius.
Fix smooth cutoffs $m_0^*,j_0:\R\to[0,1]$ with $m_0^*=1$ on
$(-\infty,0]$, $m_0^*=0$ on $[1,\infty)$, $j_0=0$ on
$(-\infty,0]$, and $j_0=1$ on $[1,\infty)$.  Define
\begin{equation}\label{four:maximal:def:penalty}
 \mathcal P=C_N(\rho_0+v\rho_1),\qquad
 m_0(t)=m_0^*(N(t+\eps)),\qquad
 \Xi=\delta_0\cT+\rho+\delta_0\eta_N a\sigma-m_0(t)\mathcal P.
\end{equation}
Proposition~\ref{four:maximal:prop:floor} below fixes a single $\delta_0>0$ independent
of $N$ and then a sufficiently large polynomial $C_N$.
Choose a smooth $N_B>1+\sup_B|H|$ on $B$ and write
\begin{equation}\label{four:maximal:def:boundary-datum}
 \mathcal B=-H+w_\nu(F-P_B)-N_Bd,\qquad P_B=\tr_BK.
\end{equation}
The system whose solution we seek is
\begin{equation}\label{four:maximal:def:system}
 \begin{cases}
  \tr_{\Achi}(K+H^f)=F=h+C,&\text{in }\Omega_R,\\
  \divg V=u\Xi,&\text{in }\Omega_R,\\
  h=b_*,\quad V_\nu/u=\mathcal B,&\text{on }B,\\
  f=Z=0,&\text{on }S_R.
 \end{cases}
\end{equation}
In the boundary datum, $F$ always means the prescribed trace right side,
so the datum contains no unprescribed second derivatives.

The two geometric signs follow directly for any smooth solution of
\eqref{four:maximal:def:system}.  Maximality and Proposition~\ref{four:maximal:prop:scalar-identity}
give the exact decomposition
\begin{equation}\label{four:maximal:def:scalar-sign}
 \frac12e^{2t}R_{\ghat}
 =\bigl(\mu+J(w)+w(C)-\rho\bigr)
  +(1-\delta_0)\cT+(1-\delta_0)\eta_N a\sigma+m_0\mathcal P.
\end{equation}
The bracket is nonnegative by \eqref{four:maximal:def:slack} and $|w|\le1$, and all
other terms are nonnegative when $\delta_0<1$.  On $B$, subtracting
\eqref{four:maximal:def:boundary-datum} from \eqref{four:maximal:def:boundary-flux} gives
$H+3\partial_\nu t=-N_B$, and therefore
\begin{equation}\label{four:maximal:def:geometric-signs}
 R_{\ghat}\ge0,\qquad \widehat H=-e^{-t}\sqrt d\,N_B<0.
\end{equation}
These conclusions use the actual boundary expansion $H+P_B=b_*$.
They require no assumption that $H$ vanishes.

Here is a homotopy to a scalar system with a trivial endpoint.
Let
\[
 V_\lambda=u\Achi\bigl(3\grad Z+\lambda K(w,\cdot)\bigr),
 \qquad 0\le\lambda\le1.
\]
In the first three stages impose $\divg V_\lambda=u\Xi$ and
\begin{equation}\label{four:maximal:def:homotopy-boundary}
 \frac{(V_\lambda)_\nu}{u}
 =\bigl[1-(1-\lambda)j_0(t)\bigr]
   \bigl[\mathcal B-(1-\lambda)
                (\Achi K(w,\cdot))_\nu\bigr].
\end{equation}
All outer values remain $f=Z=0$.  Let $\nu_s=\grad s/|ds|$ in the fixed
collars.  Choose smooth bounded functions $D_0(x,w)$ and $D_1(x,w)$,
supported in those collars and with uniformly bounded derivatives for
$|w|\le1$, as follows:
\begin{itemize}
 \item $D_0\le0$, $D_0=0$ when $w\cdot \nu_s\ge0$, and
 $D_0(x,-\nu)<-2\sup_B|H|$ on $B$;
 \item $D_1=A_1\zeta_B(x)w\cdot \nu_s$, where $\zeta_B=1$ near $B$ is a
 fixed collar cutoff and $A_1>1+\sup_B|H|$.
\end{itemize}
Such $D_0$ is obtained by a fixed smooth function of $w\cdot \nu_s$, zero
on the nonnegative half-line and sufficiently negative at $-1$, times
a collar cutoff.  In particular $D_0(x,0)=D_1(x,0)=0$.
The stages, with their endpoints included, are:
\begin{enumerate}[label=\arabic*.]
 \item Decrease $\lambda$ from one to zero, keeping $F=h+C$ and
 $h|_B=b_*$.
 \item Keep $\lambda=0$ and increase $\alpha$ from zero to one, with
 $F=h+C+\alpha D_0$ and $h|_B=b_*$.
 \item Keep $\lambda=0$ and increase $\zeta$ from zero to one, with
 \begin{equation}\label{four:maximal:def:stage-three}
  F=h+(1-\zeta)(C+D_0)
       +\zeta\bigl(\tr_{\Achi}K+D_1\bigr),
  \qquad h|_B=(1-\zeta)b_*.
 \end{equation}
 \item Keep the terminal trace problem
 \begin{equation}\label{four:maximal:def:terminal-trace}
  \tr_{\Achi}H^f=h+D_1,\qquad f|_{B\cup S_R}=0,
 \end{equation}
 and multiply both $\Xi$ and the right side of
 \eqref{four:maximal:def:homotopy-boundary} by a common factor $\eta$, decreasing
 from one to zero.
\end{enumerate}
The prescriptions agree at every junction.  In every stage the
derivative of $F$ with respect to $h$ is one.  The form $\cT$ is always
computed from the actual trace equation and the actual tensor $K+H^f$,
including when $\lambda=0$.

Two features of this choice will be used in the estimates.
First, in stage three the limiting values at $w=\nu,-\nu$, with
$\chi=0$ and the assigned face height, obey
\begin{equation}\label{four:maximal:def:face-direction-signs}
 F(x,\nu)\ge P_B+H,\qquad F(x,-\nu)\le P_B-H.
\end{equation}
At $\zeta=0$ the first equality follows from $b_*=P_B+H$ and
$D_0(x,\nu)=C|_B=0$; the second follows from the choice of $D_0$.
At $\zeta=1$ the two values are $P_B\pm A_1$.  Convex combination proves
\eqref{four:maximal:def:face-direction-signs} for intermediate $\zeta$.
Second, the terminal trace equation forces $f=0$ for every
$C^{1,\alpha}$ trial $Z$ and its classical trace solution: at a positive interior maximum the left side is nonpositive and
the right side is $\eta_N f>0$, while at a negative interior minimum the
signs reverse.  Boundary values are zero.  Thus in stage four
\begin{equation}\label{four:maximal:def:terminal-reduction}
 f=0,\quad t=Z,\quad d=\chi=1,\quad u=L_0(t),\quad
 \cT=\tfrac12|K|^2+3(p+1)|dt|^2.
\end{equation}

\begin{lemma}[Height and initial end bounds]\label{four:maximal:lem:height-end}
Every smooth homotopy solution satisfies $|h|+|F|\le C_0$ for a fixed
constant independent of $N,R$ and the homotopy parameter.  There are
fixed $r_0,C_2$ such that, for all sufficiently large $R$ depending on
$N,\eps$, the first three stages satisfy
\begin{equation}\label{four:maximal:def:height-end}
 |h|\le C_2(r^{-\gamma}-R^{-\gamma})\quad(r_0\le r\le R),
 \qquad |\grad f|\big|_{S_R}\le C_3\eta_N^{-1}R^{-1-\gamma}.
\end{equation}
Moreover $t>-\eps$ on $S_R$.  None of these assertions assumes an
interior lower bound for $t$.
\end{lemma}

\begin{proof}
At an interior extremum of $h$, we have $w=0$, $\Achi=\Id$,
$D_0=D_1=0$, and $\tr K=0$.  The trace equation and its positive
Hessian coefficient show that a minimum has $h\ge0$ and that a maximum
has $h\le\|C\|_\infty$.  The boundary heights lie in $[b_*,0]$.
Consequently
\begin{equation}\label{four:maximal:def:height-range}
 b_*\le h\le\|C\|_\infty
\end{equation}
in stages one through three; stage four has $f=0$.
All remaining terms in the trace prescriptions are uniformly bounded
for $|w|\le1$, $0\le\chi\le1$, so $F$ has a fixed bound as well.

Outside the compact supports the trace equation takes the form
\[
 \beta\tr_{\Achi}\Hess h+s_K\tr_{\Achi}K=h,
 \qquad \beta=\frac{l}{\eta_N\sqrt D},\qquad 0\le s_K\le1.
\]
Here $s_K=1$ in the first two stages and $s_K=1-\zeta$ in the third.
At a nonzero test gradient $\beta=a/|\grad h|$.
Test with $\psi=C_2(r^{-\gamma}-R^{-\gamma})$.  Since the test axis is
radial, the end metric estimates give, uniformly for $0<\chi\le1$,
\[
 \frac{\tr_{\Achi}\Hess\psi}{|\grad\psi|}
 =\frac{-3+(\gamma+1)\chi}{r}+O(r^{-1-q})
 \le-\frac{c_1}{r},\qquad c_1=(2-\gamma)/2>0,
\]
after fixing $r_0$ sufficiently large.  In addition, maximality gives
\begin{equation}\label{four:maximal:def:K-trace-estimate}
 |\tr_{\Achi}K|
 \le K_0r^{-1-q}\min(1,(1+N/3)a^2).
\end{equation}
Indeed $\tr_{\Achi}K=(\chi-1)K(e_f,e_f)$ and
$1-\chi=(3+p)v/(3+pv)\le\min(1,(1+N/3)v)$.
These estimates hold for every value of $t$.

Put $c_\gamma=1-2^{-\gamma}$ and
$H_0=\max(|b_*|,\|C\|_\infty)$.  Choose the fixed coefficient $C_2$ so that
\[
 C_2\ge\max\left\{1,
 \frac{2H_0r_0^\gamma}{c_\gamma},
 \frac{2K_0r_0^{\gamma-1-q}}{c_\gamma}\right\}.
\]
For $R\ge2r_0$, this dominates the heights at $r=r_0$.
On $r\le R/2$, the tensor term in \eqref{four:maximal:def:K-trace-estimate} is at
most $\psi/2$, since $\gamma<1+q$.  On $r\ge R/2$, use
$\min(1,(1+N/3)a^2)\le\sqrt{1+N/3}\,a$ and require
\[
 R\ge2\left(\frac{2K_0\sqrt{1+N/3}}{c_1}\right)^{1/q}.
\]
The tensor error is then at most $c_1a/(2r)$.

If $h-\psi$ had a positive maximum, it would occur in the interior of
this annulus.  The gradients there agree and the Hessian of $h$ is no
larger than the Hessian of $\psi$.  The coefficients use the actual $Z$
and this common gradient, so they agree at the contact regardless of
the two heights.  On the first half of the annulus the equation gives
$h\le-c_1a/r+\psi/2<\psi$; on the second half it gives
$h\le-c_1a/(2r)<0$.  Both are contradictions.
At a negative minimum of $h+\psi$ the Hessian inequality reverses,
while its axis and $\chi$ are unchanged on reversing the gradient;
the same estimates give the lower comparison.  This proves the height
bound in \eqref{four:maximal:def:height-end}.

Because $h=0$ at $S_R$, the two comparisons bound its normal derivative
there by $C_3R^{-1-\gamma}$, with $C_3=2C_2\gamma$ after increasing $r_0$
so that $|dr|\le2$.  Tangential derivatives vanish, giving the asserted
bound for $|\grad f|$.  Finally, at $S_R$ the strictly increasing function
$G(t)=t+\frac16\log(1+l(t)^2\sigma^2)$ has $G(t)=Z=0$.
Since $l(-\eps)=\ell$, its root exceeds $-\eps$ whenever
$\ell^2\sigma^2<e^{6\eps}-1$.  It suffices to require
\begin{equation}\label{four:maximal:def:outer-radius-floor}
 R^{2\gamma+2}>
 \frac{C_3^2e^{\eps N}}{e^{6\eps}-1}.
\end{equation}
In stage four the conclusion follows instead from $f=0$ and $Z=0$ on
$S_R$.  Every restriction on $R$ made here depends only on the fixed
data, $N,\eps$, and is compatible with $R\to\infty$ at fixed $N$.
\end{proof}

\subsection{Exclusion of the lower boundary of the admissible range}

The next argument concerns any smooth homotopy solution whose minimum
value of $t$ is $-\eps$.  It supplies the strict exclusion needed for
continuation before any global gradient or Hessian estimate is available.

\begin{lemma}[Gauss comparison at a minimum]\label{four:maximal:def:gauss-minimum}
At an interior minimum of $t$,
\begin{equation}\label{four:maximal:def:graph-gauss}
 \tfrac12e^{2t}R_{\ghat}
 \le\tfrac12R_g-v\Ric_g(e_f,e_f)+\mathcal S(H^f),
\end{equation}
where, for a symmetric tensor $A$,
\begin{equation}\label{four:maximal:def:gauss-form}
 \mathcal S(A)=\tfrac13(\tr_\perp A)^2
 -\tfrac12|A_\perp^0|^2
 +dA_{ee}\tr_\perp A-d|A_{e\perp}|^2.
\end{equation}
At every point where $dt=0$, there is an absolute $c_2>0$ such that
\begin{equation}\label{four:maximal:def:floor-comparison}
 \cT-\mathcal S(H^f)
 \ge c_2\cT-\Pi_N\bigl(vF^2+|K|^2\bigr).
\end{equation}
At $df=0$ the terms multiplied by $v$ vanish and any axis may be used.
\end{lemma}

\begin{proof}
Lemma~\ref{four:sys:gauss-lemma}, with $e=e_f$, uses exactly
\eqref{four:maximal:def:gauss-form}.  At a minimum $\Hess t\ge0$,
so its stationary Gauss identity gives
\eqref{four:maximal:def:graph-gauss}.

The two algebraic identities used by its floor estimate specialize to
\begin{align}\label{four:maximal:def:gauss-difference}
 \cT-\mathcal S(S_f)
 ={}&|T^0|^2+(1+d)|M|^2\notag\\
 &+\chi(1+d-2\chi/3)j^2+(\chi/3-d)Fj+F^2/3
\end{align}
and
\begin{equation}\label{four:maximal:def:d-chi-ratio}
 \frac{d^2}{\chi}=\frac{d(3+pv)}3\le\frac{N+3}3.
\end{equation}
They are \eqref{four:sys:stationary-difference} and the first ratio
in \eqref{four:sys:polynomial-ratios}.  The common proof splits at
$(1+p)v$ small, uses stationary coercivity, and replaces $S_f$ by
$H^f=S_f-K$.  It yields an absolute positive $c_2$ and a polynomial
$\Pi_N$ in precisely \eqref{four:maximal:def:floor-comparison},
including the smooth extension to $df=0$.
\end{proof}

\begin{lemma}[Coefficient derivatives]\label{four:maximal:def:coefficient-lemma}
The following estimates hold along every smooth trace solution:
\begin{equation}\label{four:maximal:def:derivative-bounds}
 |\nabla w|+|\nabla\Achi|+|d\log u|_{\Achi}
 \le\Pi_N(\sqrt{\cT}+|K|).
\end{equation}
Consequently, for each fixed $\eta_1>0$,
\begin{equation}\label{four:maximal:def:drift-divergence}
 \frac{|\divg(u\Achi K(w,\cdot))|}{u}
 \le\eta_1\cT+\Pi_N\bigl(|K|^2+a|\nabla K|\bigr).
\end{equation}
In each of the first three trace stages,
\begin{equation}\label{four:maximal:def:F-derivative}
 w(F)\ge\eta_N a\sigma-\eta_1\cT
   -\Pi_N\bigl(\mathbf1_{\mathcal K}+|K|^2+a|\nabla K|\bigr),
\end{equation}
where $\mathcal K$ is a fixed compact enlargement of the supports of
$C,D_0,D_1$.  For fixed $\eta_1$ the bounds remain polynomial in $N$.
\end{lemma}

\begin{proof}
Differentiating $w=l\grad f/\sqrt D$ covariantly gives
\begin{equation}\label{four:maximal:def:w-derivative-exact}
 \nabla_iw=\Ad(H^f_{i\cdot})+pd\,w\,\partial_i t,
 \qquad
 d\log u=(p+pv+2)dt+H^f(w,\cdot).
\end{equation}
In the first formula, $d$ is the scalar $D^{-1}$.  Relative to the axis,
the $\Ad$ factor multiplies the axial output of $H^f$ by $d$.
The transverse and mixed entries are controlled by
\eqref{four:maximal:def:coercivity}; the possibly large axial entry occurs as
$d(j-K_{ee})$ and is also controlled because $d\le\sqrt d$.
Furthermore $d|dt|\le|dt|_{\Ad}$.  This proves the estimate for
$\nabla w$ without an inverse power of $l$ or $d$.
For the second formula use $\chi\le d$: the weighted norm of its first
term is at most $(2N+2)|dt|_{\Ad}$, while its last term is bounded by
$a(|M-K_{e\perp}|+\sqrt\chi\,|j-K_{ee}|)$.

Write $\Achi=\Id-c(p,v)w\otimes w$, where
$c(p,v)=(3+p)/(3+pv)$.  The relevant derivatives are
\[
 \partial_pc=\frac{3d}{(3+pv)^2},\qquad
 \partial_vc=-\frac{p(3+p)}{(3+pv)^2},\qquad
 |\nabla v|\le2a|\nabla w|.
\]
The $p'$ contribution retains its factor $d$, and is bounded by a
polynomial times $d|dt|$.  All other terms are bounded by a polynomial
times $|\nabla w|$.  This proves \eqref{four:maximal:def:derivative-bounds}.

In \eqref{four:maximal:def:drift-divergence}, differentiating $u$ contributes at most
\[
 |d\log u|_{\Achi}|K(w,\cdot)|_{\Achi}.
\]
The other derivatives contribute at most
\[
 |K|(|\nabla\Achi|+|\nabla w|)+a|\nabla K|.
\]
Apply \eqref{four:maximal:def:derivative-bounds} and then Young's inequality to
$\Pi_N|K|\sqrt\cT$.  Squaring a polynomial remains polynomial, proving
\eqref{four:maximal:def:drift-divergence} with the asserted dependence.

For \eqref{four:maximal:def:F-derivative}, the common height term gives exactly
$w(h)=\eta_N a\sigma$.  Derivatives of $C,D_0,D_1$ cost, on $\mathcal K$,
a fixed multiple of $1+|\nabla w|$, hence an arbitrarily small fixed
multiple of $\cT$ plus a polynomial compact error and $\Pi_N|K|^2$.
The only other term is $w(\tr_{\Achi}K)$ in stage three; it is bounded
in absolute value by a fixed multiple of
$a|\nabla K|+a|K||\nabla\Achi|$, handled in the same way.
All stage parameters lie in $[0,1]$, so the estimates are uniform in
them.  This proves the last assertion.
\end{proof}

\begin{proposition}[Floor exclusion]\label{four:maximal:prop:floor}
There is a choice $0<\delta_0<1/4$ independent of $N$ and a positive
polynomial $C_N$ such that, for every sufficiently large allowed $R$,
no smooth solution of any of the four homotopy stages has
\[
 \min_{\overline{\Omega_R}}t=-\eps.
\]
These choices use the fixed data and the preceding pointwise and height
estimates only.  In particular they do not use existence, an upper bound
for $Z$, or a uniform Hessian estimate.
\end{proposition}

\begin{proof}
We first record all end losses that will occur.  There is a fixed
constant $C_*$ with
\begin{equation}\label{four:maximal:def:floor-weight-domination}
 \mathbf1_{\mathcal K}+|K|^2+vF^2+v|\Ric_g|+a|\nabla K|
 \le C_*(\rho_0+v\rho_1).
\end{equation}
On a fixed compact set this follows from positivity of $\rho_0$ and
Lemma~\ref{four:maximal:lem:height-end}.  Outside that set,
$|F|\le |h|+|K|$ in all three trace stages, and the same lemma controls
$h^2$ by a fixed multiple of $r^{-2\gamma}$.
The metric and tensor decay give
$|\Ric_g|=O(r^{-2-q})$, $|K|^2=O(r^{-2-2q})$ and
$|\nabla K|=O(r^{-2-q})$.  Since $2\gamma<2+q$ and $2+2q>4+\delta$,
the curvature and squared tensor terms have the stated bounds.
For the derivative term the precise estimate is
\[
 ar^{-2-q}
 \le\tfrac12a^2r^{-2\gamma}
       +\tfrac12r^{-4-2q+2\gamma}
 \le\tfrac12v\rho_1+\tfrac12\rho_0,
\]
where $2(q-\gamma)>\delta$ was used.  This proves
\eqref{four:maximal:def:floor-weight-domination}.  Also $\rho\le C_\rho\rho_0$ for a
fixed $C_\rho$, by its prescribed tail and compact positivity of $\rho_0$.

Suppose first that the minimum occurs in the interior in one of the
first three stages.  Then $dt=0$ and $\Hess t\ge0$ there.
Proposition~\ref{four:maximal:prop:scalar-identity} and
\eqref{four:maximal:def:graph-gauss} imply, using $\tr K=0$,
\[
 u^{-1}\divg V
 \ge\mu-\tfrac12R_g+J(w)+v\Ric_g(e_f,e_f)
       +\cT-\mathcal S(H^f)+w(F).
\]
Here $\mu-R_g/2=-|K|^2/2$ and
$|J(w)|\le C_*a|\nabla K|$.
Use \eqref{four:maximal:def:floor-comparison}, \eqref{four:maximal:def:F-derivative}, and
\eqref{four:maximal:def:floor-weight-domination}; then pass from $V$ to
$V_\lambda=V-(1-\lambda)u\Achi K(w,\cdot)$ by
\eqref{four:maximal:def:drift-divergence}.  Choose the two fixed Young errors smaller
than $c_2/4$.  This yields an absolute $c_3>0$ and a polynomial bound
$\Pi_N^{\mathrm f}$ such that at the proposed minimum
\begin{equation}\label{four:maximal:def:floor-divergence}
 u^{-1}\divg V_\lambda
 \ge c_3\cT+\eta_N a\sigma
       -\Pi_N^{\mathrm f}(\rho_0+v\rho_1).
\end{equation}
In particular $c_3$ is independent of $N$.

Fix $\delta_0<\min(c_3/2,1/4)$.  At the floor $m_0=1$, and the prescribed
right side in the first three stages is
$\delta_0\cT+\rho+\delta_0\eta_N a\sigma-C_N(\rho_0+v\rho_1)$.
The lower bound in \eqref{four:maximal:def:floor-divergence} exceeds it by at least
\[
 (c_3-\delta_0)\cT+(1-\delta_0)\eta_N a\sigma
 +(C_N-\Pi_N^{\mathrm f})(\rho_0+v\rho_1)-\rho.
\]
Choose $C_N\ge\Pi_N^{\mathrm f}+C_\rho+1$.  The last display is strictly
positive because $\rho_0>0$, contradicting the equation.  This choice is
polynomial and independent of $R$ and of the solution.

In stage four use \eqref{four:maximal:def:terminal-reduction}.  At an interior floor
minimum the left side is
$u^{-1}\divg(3u\grad t)=3\Delta t\ge0$,
while the right side is
\[
 \eta\bigl(\delta_0|K|^2/2+\rho-C_N\rho_0\bigr).
\]
Increasing $C_N$ by a fixed constant makes the bracket strictly negative
everywhere, since $|K|^2\le C_*\rho_0$.  Thus an interior floor minimum
is impossible whenever $\eta>0$.

An inner boundary minimum in the first three stages is also impossible.
At $t=-\eps$ we have $j_0=0$.  Since
$V_\lambda=V-(1-\lambda)u\Achi K(w,\cdot)$, the drift correction in
\eqref{four:maximal:def:homotopy-boundary} cancels, leaving $V_\nu/u=\mathcal B$.
Lemma~\ref{four:maximal:lem:boundary-identity} therefore gives
\[
 3\partial_\nu t=-H-N_B<0.
\]
This contradicts the nonnegative derivative into $\Omega_R$ required
at a boundary minimum.  In stage four, where $f=0$, the corresponding
formula is $3\partial_\nu t=\eta(-H-N_B)<0$ for $\eta>0$.
The outer boundary is excluded by Lemma~\ref{four:maximal:lem:height-end}.

It remains to consider $\eta=0$.  The last-stage equation and boundary
conditions then reduce to
\[
 \divg(3L_0(Z)\grad Z)=0,\qquad
 \partial_\nu Z|_B=0,\qquad Z|_{S_R}=0.
\]
Multiplying by $Z$ and integrating gives
$\int_{\Omega_R}3L_0(Z)|\grad Z|^2\dd V_g=0$.
Thus $Z=t=0$, which is strictly above the floor.  This treats every
stage and every possible location of a floor minimum.
\end{proof}

The order of the choices is now fixed: strict background data and
collars, $\eps$ and the cutoffs, the absolute small $\delta_0$, the
polynomial $C_N$, and finally sufficiently large $N$ and allowed $R$.
The estimates in the next Section concern arbitrary smooth admissible
solutions with $t\ge-\eps$.  Proposition~\ref{four:maximal:prop:floor} separately
ensures that a homotopy solution cannot meet the boundary of this range.

\section{Global a priori bounds}\label{four:maximal:sec:bounds}

Throughout this section the strict background data, collars, and
$\eps>0$ are fixed. We consider arbitrary smooth solutions of the
homotopy on $\Omega_R$ satisfying $t\ge-\eps$. Existence is not assumed:
each assertion is an estimate for every solution that might occur.
The constants denoted by $\Pi_N$ are bounded by a polynomial in $N$,
independently of the solution, homotopy parameter, and allowed radius.
Constants denoted by $C(N)$ may depend arbitrarily on the fixed integer
$N$. All integrals without an indicated graph measure use $g$.

\subsection{Separation from the threshold region}

\begin{lemma}[Compact separation]\label{four:maximal:lem:separation}
Let $Q_0$ be a compact subset of $\overline\Omega$ disjoint from
$\mathcal D_c$. There are $\eta_2>0$, $N_0$, and allowed radii
$R_{\min}(N)\to\infty$ such that, in stages~1 and~2,
\begin{equation}\label{four:maximal:bounds:separation}
 h\ge c+\eta_2\quad\hbox{on }Q_0
 \qquad(N\ge N_0,\ R\ge R_{\min}(N)).
\end{equation}
The constants are uniform over both stages.
\end{lemma}

\begin{proof}
We first identify the geometric inequality satisfied by a relaxed
limit. Take any sequence of admissible solutions with $N_i\to\infty$
and allowed $R_i\to\infty$, and write
\[
 \underline h(x)=\lim_{j\to\infty}
 \inf\{h_i(y):i\ge j,\ \dist(x,y)<j^{-1}\}.
\]
On compact sets this is a finite lower semicontinuous function, by
Lemma~\ref{four:maximal:lem:height-end}. Let a smooth function $\phi$ touch
$\underline h$ from below at an interior point $x$, with
$\grad\phi(x)\ne0$. Subtracting a small multiple of the fourth power
of the distance makes the contact strict without changing its two-jet.
Minimization on a small closed ball then gives a subsequence of local
lower contacts of $\phi$ plus constants with $h_i$, at points $x_i\to x$,
whose contact values converge to $\underline h(x)$.

At these contacts, $\eta_{N_i}\grad f_i=\grad\phi$, so
\[
 l_i\sigma_i\ge\frac{\ell_i}{\eta_{N_i}}|\grad\phi(x_i)|\longrightarrow\infty,
 \qquad d_i,\chi_i\longrightarrow0,\quad a_i\longrightarrow1.
\]
In particular this conclusion uses the floor and the fixed test
gradient, rather than an estimate for the solution gradients.
Put $\nu_\phi=\grad\phi/|\grad\phi|$. Positivity of $\Achi$ and the
Hessian inequality at a lower contact give
\[
 \tr_{A_{\chi_i}}K+
 \frac{a_i}{|\grad\phi|}
       \tr_{A_{\chi_i}}\Hess\phi
 \le F_i\le h_i.
\]
The last inequality holds throughout stages~1 and~2 because
$C\le0$ and $D_0\le0$. Passing to the limit proves
\begin{equation}\label{four:maximal:bounds:lower-test}
 \frac{\tr_{\nu_\phi^\perp}\Hess\phi}{|\grad\phi|}
       +\tr_{\nu_\phi^\perp}K\le\underline h(x).
\end{equation}
The left side is the outward expansion of the regular level of $\phi$.
Neither convergence of the homotopy parameters nor continuity bounds
for the solution coefficients were needed.

Fix $c<k<k'<0$.  Apply Lemma~\ref{n3:apriori:sublevel-support}
on $U=\operatorname{int}\Omega$ with $u=\underline h$, $Q=K$ and
$q_0=k'$.  The relaxed limit is finite and lower semicontinuous.
At a nonzero-gradient lower test whose contact value is less than $k'$,
\eqref{four:maximal:bounds:lower-test} bounds its expansion by that
value, hence by $k'$.  The lemma therefore shows that every smooth
exterior support to the closed sublevel $\{\underline h\le k\}$
at an interior frontier point has expansion at most $k'$.  In particular,
this conclusion also applies when the sublevel has a plateau at height $k$.

Lemma~\ref{four:maximal:lem:height-end} makes this sublevel bounded: for a large
fixed $r$, its lower height bound exceeds the fixed negative number
$k$. Adjoin $\mathcal D_c$ to obtain a compact closed set. At a frontier
point on $B$, each exterior support also supports $\mathcal D_c$;
tangency comparison with its smooth boundary gives expansion at most
$c$. Lemma~\ref{four:maximal:lem:support-enclosure} therefore places the union in
$\mathcal D_{c'}$ for every $k'<c'<0$.

Hausdorff right continuity at $c$ and the positive distance of $Q_0$
from $\mathcal D_c$ allow us to choose $c<k<k'<c'<0$ so that
$\mathcal D_{c'}\cap Q_0=\varnothing$. If
\eqref{four:maximal:bounds:separation} failed, choose $N_i\to\infty$, allowed
$R_i\ge i$, and $x_i\in Q_0$ with $h_i(x_i)\le c+i^{-1}$.
A subsequence has $x_i\to x\in Q_0$, and then
$\underline h(x)\le c<k$, contradicting that containment.
This also proves uniformity over the two stages. Enlarging the allowed
radius function never invalidates the assertion.
\end{proof}

\subsection{Collar barriers and polynomial trace estimates}

\begin{lemma}[Inner collar slopes]\label{four:maximal:lem:collar-slopes}
For all sufficiently large $N$ and allowed $R$, there are fixed
constants $0<c_5<C_6$ such that
\begin{equation}\label{four:maximal:bounds:collar-slopes}
 \frac{c_5}{\eta_N}\le\partial_\nu f\le\frac{C_6}{\eta_N}
 \quad\hbox{in stages 1 and 2},\qquad
 |\partial_\nu f|\le\frac{C_6}{\eta_N}
 \quad\hbox{in stage 3}.
\end{equation}
Here $\nu$ points into $\Omega$. In stage~4, $f=0$.
\end{lemma}

\begin{proof}
Use a fixed short collar $0\le s\le s_0$ from
Lemma~\ref{four:maximal:lem:threshold-collars}, with $C=-\kappa s$,
$\nu_s=\grad s/|\dd s|$, and $0<c_s\le|\dd s|\le C_s$.
All its geometric derivatives are bounded, and
$P_s+H_s\ge c$, where $P_s=\tr_{\{s=\mathrm{constant}\}}K$.
The finitely many outer collar faces satisfy
Lemma~\ref{four:maximal:lem:separation}.

For a positive-slope test $h_0=c+\psi(s)$, the trace operator, with
$t$ computed from the solution's $Z$ and the common contact gradient,
is exactly
\begin{equation}\label{four:maximal:bounds:collar-operator}
 P_s+aH_s+\chi K(\nu_s,\nu_s)
 +a\chi\left(
    \frac{\Hess s(\nu_s,\nu_s)}{|\dd s|}
       +|\dd s|\frac{\psi''}{\psi'}\right).
\end{equation}
For a test with negative slope, replace $a$ by $-a$ in this formula.
If $|\psi'|$ lies between fixed positive constants, the floor gives
\begin{equation}\label{four:maximal:bounds:steep-coefficients}
 d\le C(\eta_N/\ell)^2=Ce^{-\eps N},\qquad
 1-a\le d,\qquad
 \frac{N\chi}{d}=\frac{3N}{3+pv}\ge\frac{3N}{3+N}.
\end{equation}
Thus $a\ge1/2$ and $N\chi\ge c_4d$ for large $N$.

Choose a fixed smooth $\vartheta_1:[0,\infty)\to[0,1]$ equal to one
on $[0,1]$ and zero on $[2,\infty)$, and put
$b_N(s)=\Lambda_* N\vartheta_1(Ns)$, where $\Lambda_*>0$ is a fixed collar
constant, independent of $N$. The constant will be
chosen large; its crucial property is
\[
 0\le\int_0^{s_0}b_N(s)\dd s\le2\Lambda_*.
\]
For the lower barrier in stages~1 and~2, prescribe
$\psi'(s)=q_-\exp(\int_0^s b_N)$ and $\psi(0)=0$.
For the chosen value of $\Lambda_*$, take $q_->0$ so small that $\psi'\le\kappa/2$ and
$\psi(s_0)<\eta_2$. Positive normal direction makes $D_0=0$.
Subtracting the prescribed trace right side from
\eqref{four:maximal:bounds:collar-operator} gives a residual at least
\begin{equation}\label{four:maximal:bounds:lower-residual}
 \kappa s/2-Cd+a\chi|\dd s|b_N(s).
\end{equation}
The constant $C$ in this residual depends only on collar geometry.
We choose $\Lambda_*$ so large that the last term dominates $2Cd$ on
$[0,1/N]$. On $[1/N,s_0]$, the first term dominates $Cd$ for large
$N$, because $Ne^{-\eps N}\to0$, and the last term has a favorable
sign. This is a strict lower barrier. Its outer value is below $h$
by separation and its inner value agrees with $h$.

For the upper barrier take
$\psi'(s)=q_+\exp(-\int_0^s b_N)$.
Choose its fixed minimum slope larger than twice the bounded linear
variation of $P_s+H_s-c+\kappa s$, and choose $q_+$ still larger if
needed to dominate the height bound at $s_0$. The residual is now at
most $-c_7s+Cd-a\chi|\dd s|b_N$, for fixed $c_7>0$.
The same two intervals prove strict negativity. The bounded integral
of $b_N$ leaves all lower and upper test slopes between fixed positive
constants, independently of $N$.

For completeness, at an interior positive maximum of
$h_0-h$, the test Hessian is at most the solution Hessian, whereas its
height is larger. The positive principal matrix and the strictly
increasing dependence of $F$ on $h$ contradict a positive test residual.
This proves lower comparison; reversing the difference proves upper
comparison. In both comparisons $Z$ and the gradient are the same
at contact, so all coefficients other than the height agree.
Taking inward derivatives at $s=0$ gives the first assertion of
\eqref{four:maximal:bounds:collar-slopes}.

In stage~3 the face height is $b_\zeta=(1-\zeta)c$.
Use $b_\zeta\pm\psi(s)$ with fixed large positive slopes and
$\psi''/\psi'=-b_N$. At $s=0$ the two unit-direction inequalities
in \eqref{four:maximal:def:face-direction-signs} say that the positive-slope residual is nonpositive
and the negative-slope residual is nonnegative. Replacing $\pm \nu_s$ by
$\pm a\nu_s$ costs $O(d)$, including the term $\chi K(\nu_s,\nu_s)$; moving to
distance $s$ costs $O(s)$, uniformly in $\zeta$. The change of height
then supplies a strict linear margin, with the desired sign, when the
minimum test slope exceeds the fixed spatial error constant.
The term involving $b_N$ has the correct sign for both tests and
dominates the $O(d)$ error on $[0,1/N]$. On the remainder the linear
margin dominates it. The height bound handles $s=s_0$.
These two comparisons bound the slope in absolute value in stage~3.
The terminal trace equation has $f=0$, as established in
Section~\ref{four:maximal:sec:deformation}.
\end{proof}

\begin{lemma}[Weighted trace and boundary flux]\label{four:maximal:lem:weighted-trace}
Let $\mathcal C$ be a fixed union of collars. In stages~1 and~2, every
nonnegative $C^1$ function $\varphi$ satisfies
\begin{align}
 \int_B L_0\varphi
 &\le\Pi_N\int_{\mathcal C}
  u\bigl((1+\sqrt{\cT})\varphi+|\dd\varphi|_{\Achi}\bigr),
       \label{four:maximal:bounds:trace-weighted}\\
 \int_B\varphi
 &\le\Pi_N\int_{\mathcal C}
  \sqrt D\bigl((1+\sqrt{\cT})\varphi+|\dd\varphi|_{\Achi}\bigr).
       \label{four:maximal:bounds:trace-unweighted}
\end{align}
The integrands vanish outside the support of $\varphi$ and its
derivative; thus both statements localize to any base patch.
Moreover
\begin{equation}\label{four:maximal:bounds:boundary-flux}
 |(V_\lambda)_\nu|\le C_7ud=C_7L_0\sqrt d
                     \le C_8(\eta_N/\ell)L_0.
\end{equation}
Consequently,
\begin{equation}\label{four:maximal:bounds:boundary-absorption}
 \int_B|(V_\lambda)_\nu|
 \le \Pi_N\frac{\eta_N}{\ell}
             \int_{\mathcal C}u(\rho+\delta_0\cT).
\end{equation}
All constants in these assertions have the stated polynomial control.
\end{lemma}

\begin{proof}
Set $W=l\grad_{\gbar}f=\sqrt d\,w$. For each covector $\xi$,
\[
 |\xi(W)|\le|\xi|_{\Ad}
       \le\sqrt{(3+N)/3}\,|\xi|_{\Achi},\qquad uW=L_0w.
\]
The second inequality follows from $d/\chi=(3+pv)/3$.
Lemma~\ref{four:maximal:lem:collar-slopes} gives $w\cdot\nu=a\ge1/2$ on $B$.
Direct differentiation yields
\begin{align*}
 \divg w&=\tr_{\Ad}H^f+dp\,w(t),\\
 u^{-1}\divg(uW)
 &=\sqrt d\bigl(\tr_{\Ad}H^f+(dp+p+2)w(t)\bigr),\\
 D^{-1/2}\divg(\sqrt D W)
 &=\sqrt d\bigl(\tr_{\Ad}H^f+dp\,w(t)\bigr).
\end{align*}
These last two expressions have absolute value at most
$\Pi_N(1+\sqrt{\cT})$ on the collars. Indeed the axial tensor term is
$d^{3/2}H^f_{ee}$, controlled by the $d j^2$ term of
Lemma~\ref{four:maximal:lem:coercivity}, and
$\sqrt d\,|w(t)|\le|\dd t|_{\Ad}$; all background $K$ terms are
bounded there. No inverse power of $l$ has occurred.

Take a fixed collar cutoff $\eta$, equal to one near $B$ and zero
near the outer collar faces. Apply the divergence theorem to
$\eta\varphi uW$. The domain's outward normal on $B$ is $-\nu$,
so the boundary integral is $-\int_B L_0a\varphi$.
Taking absolute values of the three interior terms, and using
$a\ge1/2$, proves \eqref{four:maximal:bounds:trace-weighted}.
The derivative of $\eta$ costs only a fixed multiple of $\varphi$.
Using $\sqrt D W=w$ instead proves \eqref{four:maximal:bounds:trace-unweighted}.
The same computation proves the asserted localization.

On $B$, positive slope gives $w_\nu=a$, $D_0=0$, and
$F=c=P_B+H$. Thus the expression inside the boundary multiplier of
the homotopy is
\[
 (a-1)H-N_Bd-(1-\lambda)\chi aK(\nu,\nu).
\]
The multiplier lies in $[0,1]$, and $1-a\le d$, $\chi\le d$.
This proves the first inequality in \eqref{four:maximal:bounds:boundary-flux}.
The lower collar slope gives
$\sqrt d\le(l|\partial_\nu f|)^{-1}\le C\eta_N/\ell$,
proving the rest. Finally apply \eqref{four:maximal:bounds:trace-weighted} with
$\varphi=1$. On the fixed collars, $\rho$ has a positive minimum, so
$1+\sqrt{\cT}\le C(\rho+\delta_0\cT)$ with fixed $C$.
This proves \eqref{four:maximal:bounds:boundary-absorption}.
\end{proof}

\subsection{The global high-level energy bound}

The collar estimates have already made the inner flux small with
polynomial constants.  We now obtain a bounded range for the conformal
variable at each fixed parameter.  That is a separate task: its constants
may be arbitrary at fixed parameter because the mass estimate will
return to the earlier collar bounds.

We now allow arbitrary constants $C(N)$. At sufficiently high levels
$Z>Z_0(N)>0$, the unscaled source satisfies
\begin{equation}\label{four:maximal:bounds:high-source}
 \Xi\ge\delta_0\cT+\rho+\tfrac12\delta_0\eta_N a\sigma.
\end{equation}
To see this, $\mathcal P$ is bounded at fixed $N$, and on the admissible
support of $m_0$ we have
$-\eps\le t\le-\eps+1/N$. There
$D=e^{6(Z-t)}$ and
$\eta_N a\sigma=(\eta_N/l)(\sqrt D-D^{-1/2})$ tends uniformly to infinity
as $Z\to\infty$. Outside this band the penalty vanishes.
On the full range we also have $\Xi\ge\delta_0\cT-C(N)$.

Only stage~1 requires an energy argument to handle its nonzero drift.
Choose a smooth nondecreasing $k_0$ which is zero below $Z_0$, one
above $Z_0+1$, and has bounded derivative. Testing the divergence
equation with $k_0(Z)$, whose outer trace is zero, gives
\[
 \int_{\Omega_R}uk_0\Xi
       +3\int_{\Omega_R}uk_0'|\dd Z|_{\Achi}^2
 =-\int_B k_0(V_\lambda)_\nu
       -\lambda\int_{\Omega_R}uk_0'
                   \langle K(w,\cdot),\dd Z\rangle_{\Achi}.
\]
Young's inequality absorbs half the gradient term and leaves
$C\int uk_0'|K|^2$. This is at most $C(N)$: on the derivative strip,
\[
 u=l e^{3Z-t}\le e^{1+3(Z_0+1)+\eps},
\]
and $|K|^2$ is integrable since $q>1$.
By \eqref{four:maximal:bounds:boundary-flux} and \eqref{four:maximal:bounds:trace-weighted},
the boundary cost is at most
\[
 \Pi_N\frac{\eta_N}{\ell}\int_{\mathcal C}
   u\bigl((1+\sqrt{\cT})k_0+k_0'|\dd Z|_{\Achi}\bigr).
\]
The first term is absorbed by a fixed fraction of
$\int uk_0(\delta_0\cT+\rho)$ for large $N$, because
$\Pi_N\eta_N/\ell\to0$. Young's inequality absorbs the second term
into $\int uk_0'|\dd Z|_{\Achi}^2$ and a bounded derivative-strip
integral over $\mathcal C$. We obtain
\begin{equation}\label{four:maximal:bounds:global-energy}
 \int_{\Omega_R}u\bigl[
       k_0(\cT+\rho+\eta_N a\sigma)
                    +k_0'|\dd Z|_{\Achi}^2\bigr]\le C(N).
\end{equation}
This is uniform in $R$ and in the stage~1 parameter.

Let $\Gamma$ denote the ordinary graph of $f$ in $(\Omega,g)\times\R$
with product metric $g+\dd b^2$. This is a four-dimensional graph.
Its measure and gradient satisfy, at fixed $N$,
\begin{equation}\label{four:maximal:bounds:graph-comparison}
 \dd\Gamma\asymp_N\sqrt D\dd V_g,\qquad
 |\grad_\Gamma\phi|\asymp_N|\dd\phi|_{\Achi}.
\end{equation}
These comparisons use only $\ell\le l\le e$ and
$3d/(3+N)\le\chi\le d$. Also
\[
 D^{1/3}\le C(N)(1+\eta_N a\sigma),\qquad
 e^{2Z}\sqrt D=l^{-1}uD^{1/3}.
\]
For the first inequality, put $y=l\sigma$. It is immediate for
$y\le1$; for $y\ge1$ use
$\eta_N a\sigma\ge(\eta_N/(\sqrt2 l))y$ and $y^{2/3}\le y$.
On any fixed compact base set $Q$, the positive minimum of $\rho$
and \eqref{four:maximal:bounds:global-energy} control the integral of
$u(1+\eta_N a\sigma)$ above $Z_0+1$. Below that level the floor and
$D=e^{6(Z-t)}$ bound every integrand on $Q$. Consequently
\begin{equation}\label{four:maximal:bounds:graph-l2}
 \int_{\Gamma\cap\pi^{-1}(Q)} e^{2Z}\dd\Gamma\le C(N,Q).
\end{equation}
Here $\pi$ is the graph projection. The estimate is local in the base;
the global estimate from which it follows is
\eqref{four:maximal:bounds:global-energy}.

\subsection{Graph Sobolev inequality and the exponential test}

We continue with stage~1 for the next two estimates.
For a fixed compact base patch and a function $\omega$ supported over
it, away from $S_R$ but possibly meeting $B$, we claim
\begin{equation}\label{four:maximal:bounds:graph-sobolev}
 \|\omega\|_{L^4(\dd\Gamma)}^2
 \le C(N)\int_\Gamma
       \bigl(|\grad_\Gamma\omega|^2+(1+\cT)\omega^2\bigr)\dd\Gamma.
\end{equation}
Extend a compact enlargement of the base metric smoothly to a closed
manifold and use a smooth isometric Euclidean embedding
\cite[Theorem~2, p.~59]{Nash1956}. Its product with $\R$ gives a Euclidean
immersion of the graph. The contribution to its mean curvature from
the base embedding is bounded by a fixed constant: it is a trace of
the fixed second fundamental form on a four-dimensional tangent
plane. The remaining product mean curvature has magnitude
\[
 \frac{\sqrt D}{l\sqrt{1+\sigma^2}}
 \left|\tr_\perp H^f+(1+\sigma^2)^{-1}H^f_{ee}\right|
 \le C(N)(1+\sqrt{\cT}).
\]
Indeed the prefactor is bounded at fixed $N$, while
$(1+\sigma^2)^{-1}\le C(N)d$ and the axial term is controlled by
$d j^2$ in Lemma~\ref{four:maximal:lem:coercivity}.

The Euclidean submanifold Sobolev inequality in dimension four gives
\[
 \left(\int_\Gamma |\omega|^4\dd\Gamma\right)^{3/4}
 \le C(N)\int_\Gamma
    \bigl(|\omega|^2|\grad_\Gamma\omega|
               +(1+\sqrt{\cT})|\omega|^3\bigr)\dd\Gamma
           +C\int_B|\omega|^3.
\]
We have applied the $L^1$ Sobolev inequality to $|\omega|^3$;
see Michael--Simon~\cite{MichaelSimon1973} and the precise
rectifiable-submanifold statement in
\cite[Chapter~4, \S5, Theorem~5.7, p.~98]{Simon2014}.
Its hypotheses hold for this smooth multiplicity-one immersion with
locally integrable mean curvature and compactly supported test.
To justify the displayed boundary term from the interior statement,
extend each individual smooth graph across $B$ and cut off in an
exterior strip whose width tends to zero. The integral of the cutoff
gradient tends to the boundary integral, while the strip's other
integrals tend to zero. This requires smoothness of each graph, not a
uniform curvature estimate for extensions. The induced boundary
metric is $g|_B$, since $f$ is constant there.

Apply \eqref{four:maximal:bounds:trace-unweighted} to $|\omega|^3$ and then
\eqref{four:maximal:bounds:graph-comparison}. It bounds the last term by the same
two interior terms. Each such term is at most
\[
 C(N)\left(\int_\Gamma|\omega|^4\dd\Gamma\right)^{1/2}
 \left(\int_\Gamma
       (|\grad_\Gamma\omega|^2+(1+\cT)\omega^2)\dd\Gamma\right)^{1/2}
\]
by Cauchy--Schwarz. Dividing, and squaring, proves
\eqref{four:maximal:bounds:graph-sobolev}; the case $\omega=0$ is immediate.

Next let $\psi$ be a smooth compact base cutoff of the same kind.
There is $k_*(N)\ge1$ such that for every real $k\ge k_*(N)$,
\begin{equation}\label{four:maximal:bounds:exponential-energy}
 \int_\Gamma e^{2kZ}\psi^2
        (\cT+k|\grad_\Gamma Z|^2)\dd\Gamma
 \le C(N)k\int_\Gamma e^{2kZ}
               (\psi^2+|\dd\psi|_g^2)\dd\Gamma.
\end{equation}
The constant $C(N)$ here is independent of $k$. To prove this, test
the stage~1 divergence equation with
$\varphi=\psi^2e^{2kZ}/L_0$. Write
$Q=\psi^2e^{2kZ}$ and $b=\lambda K(w,\cdot)$.
After integration, cancellation of $L_0$ gives the exact identity
\begin{align*}
 &\int_{\Omega_R}\sqrt D\,Q
       \bigl(\Xi+6k|\dd Z|_{\Achi}^2\bigr)\\
 &=-\int_B(V_\lambda)_\nu Q/L_0
   -\int_{\Omega_R}\sqrt D\,e^{2kZ}
       \bigl[2\psi\langle3\dd Z+b,\dd\psi\rangle_{\Achi}
          +2k\psi^2\langle b,\dd Z\rangle_{\Achi}\bigr]\\
 &\hspace{1cm}
   +\int_{\Omega_R}\sqrt D\,Q(p+2)
                    \langle3\dd Z+b,\dd t\rangle_{\Achi}.
\end{align*}
Use $\Xi\ge\delta_0\cT-C(N)$, $|b|_{\Achi}\le|K|$, and
Lemma~\ref{four:maximal:lem:coercivity}. For example,
\[
 (p+2)|\dd t|_{\Achi}(|\dd Z|_{\Achi}+|K|)
 \le\frac{\delta_0}{8}\cT
               +C(N)(|\dd Z|_{\Achi}^2+|K|^2).
\]
Young's inequality handles the drift and cutoff products with a
fixed fraction of $k|\dd Z|_{\Achi}^2$ and a remainder bounded by
$C(N)k e^{2kZ}(\psi^2+|\dd\psi|^2)$. Choose $k_*(N)$ large enough
to absorb the preceding $C(N)|\dd Z|_{\Achi}^2$ term.

By \eqref{four:maximal:bounds:boundary-flux} the boundary contribution is at most
$C\int_B Q$ (enlarging a fixed constant since $\eta_N/\ell\le1$).
Applying \eqref{four:maximal:bounds:trace-unweighted} to $Q$ gives, only on
$\mathcal C\cap\supp\psi$,
\[
 C(N)\int\sqrt D\,e^{2kZ}
   \bigl[(1+\sqrt{\cT})\psi^2
       +2|\psi||\dd\psi|+2k\psi^2|\dd Z|_{\Achi}\bigr].
\]
The $\sqrt{\cT}$ and $k|\dd Z|_{\Achi}$ terms are absorbed by
small fixed fractions of $\delta_0\psi^2\cT$ and
$k\psi^2|\dd Z|_{\Achi}^2$, leaving at most
$C(N)k e^{2kZ}(\psi^2+|\dd\psi|^2)$.
Finally \eqref{four:maximal:bounds:graph-comparison} proves
\eqref{four:maximal:bounds:exponential-energy}. In particular, neither the trace
cutoff nor differentiation of $1/L_0$ introduces a constant exponential
in $k$.

\subsection{Moser iteration and the full range}

\begin{proposition}[Bounded variables]\label{four:maximal:prop:bounded-range}
For every sufficiently large fixed $N$ and all sufficiently large
allowed $R$, every admissible smooth homotopy solution satisfies
\begin{equation}\label{four:maximal:bounds:bounded-range}
 |f|\le C_0/\eta_N,\qquad
 -\eps\le t\le Z\le C(N),\qquad |\grad f|\le C(N).
\end{equation}
The constants are independent of $R$, the solution, and the homotopy
parameter. No uniform Hessian or coefficient continuity estimate is
used in obtaining these bounds.
\end{proposition}

\begin{proof}
We first finish stage~1. Apply \eqref{four:maximal:bounds:graph-sobolev} to
$\omega=\psi e^{kZ}$ and then use
\eqref{four:maximal:bounds:exponential-energy}. If $\psi=1$ on a smaller patch
and $|\dd\psi|\le C/\mathrm{gap}$, the resulting norm step is
\[
 \|e^Z\|_{L^{4k}(\text{smaller graph})}
 \le\bigl[C(N)k^2(1+\mathrm{gap}^{-2})\bigr]^{1/(2k)}
       \|e^Z\|_{L^{2k}(\text{larger graph})}.
\]
Choose $k_j=2^j k_*(N)$ and geometrically decreasing gaps between
fixed nested base patches. The logarithms of the factors sum because
$\sum_j (1+j)/2^j<\infty$. Writing $S(r)=\sup_{Q_r}e^Z$ for a
nested family of compact patches gives
\[
 S(r')\le C(N)(r''-r')^{-A(N)}
          \|e^Z\|_{L^{2k_*}(\Gamma\cap\pi^{-1}(Q_{r''}))}.
\]
Interpolation with \eqref{four:maximal:bounds:graph-l2} yields
\[
 \|e^Z\|_{2k_*}
     \le S(r'')^{1-1/k_*}\|e^Z\|_2^{1/k_*}.
\]
If $k_*=1$ this already proves the bound. Otherwise Young's inequality
gives, for every $\alpha>0$,
\[
 S(r')\le\alpha S(r'')
       +C(N,\alpha)(r''-r')^{-A'(N)}.
\]
Iterate over radii increasing geometrically to a fixed larger radius
and take $\alpha<2^{-A'(N)-1}$. The sum of the remainder terms is
finite. The final term $\alpha^j S(r_j)$ tends to zero for each
individual smooth solution, since all $Q_{r_j}$ lie in a fixed compact
patch. The resulting bound is independent of that individual
supremum. Finitely many such patches prove an upper bound for $Z$
on any fixed compact base set, including $B$.

To extend the bound along the end, use \eqref{four:maximal:def:drift-divergence},
with a fixed small coefficient:
\[
 u^{-1}|\divg(u\Achi K(w,\cdot))|
 \le\frac{\delta_0}{4}\cT
       +\Pi_N(|K|^2+a|\grad K|).
\]
Since $a\sigma= l\sigma^2/\sqrt D\ge e^{-1}a^2$,
Young's inequality bounds the second term by
\[
 \frac{\delta_0}{4}\eta_N a\sigma
       +C(N)(|K|^2+|\grad K|^2).
\]
The last expression, apart from its favorable term, is $o(\rho)$
on the end. Indeed $2+2q>4+\delta$ and
$|\grad K|^2=O(r^{-4-2q})$.
Beyond a sufficiently large fixed-$N$ sphere, the high-level source
\eqref{four:maximal:bounds:high-source} therefore implies
\[
 \divg(3u\Achi\grad Z)>0
 \quad\hbox{where }Z>Z_0(N).
\]
At an interior maximum its left side is
$3u\Achi:\Hess Z\le0$, a contradiction.
The compact bound controls the inner sphere and $Z=0$ on $S_R$.
This proves the global upper bound in stage~1.

In stages~2 and~3, $\lambda=0$. The same high-level source is
strictly positive, so it excludes interior high maxima directly.
The boundary slope bound implies
\[
 Z=t+\tfrac16\log(1+l^2|\partial_\nu f|^2)
       \le t+\tfrac16\log(1+e^2 C_6^2/\eta_N^2)
       \quad\hbox{on }B.
\]
Thus sufficiently high boundary $Z$ forces $t\ge1$. At such points
$j_0(t)=1$, and the boundary prescription gives zero conormal flux.
The boundary point principle excludes a nonconstant high maximum
there: in a neighborhood of the maximum the individual smooth
operator is elliptic and $\divg(3u\Achi\grad Z)>0$, so its outward
conormal derivative at a boundary maximum is strictly positive,
contrary to the zero datum. This uses ellipticity only for each
individual solution. The outer Dirichlet value is zero, so the upper
bound follows uniformly through these stages.

In stage~4, $f=0$ and $t=Z$. For positive scaling $\eta$, sufficiently
high levels lie outside the penalty band and have source
$\eta(\delta_0\cT+\rho)>0$; the same maximum and boundary point
argument applies. At zero scaling, multiplying the homogeneous
equation by $Z$ gives $\int 3L_0|\dd Z|^2=0$, and the outer value yields
$Z=0$. These bounds are uniform even as $\eta\downarrow0$.

Finally the height bound is Lemma~\ref{four:maximal:lem:height-end}, and
$t\le Z$ follows from $D\ge1$. The implicit equation gives
\[
 D=e^{6(Z-t)}\le e^{6(C(N)+\eps)},\qquad
 |\grad f|^2=(D-1)/l^2\le\ell^{-2}e^{6(C(N)+\eps)}.
\]
This proves \eqref{four:maximal:bounds:bounded-range}.
\end{proof}

\begin{remark}\label{four:maximal:bounds:parameter-order}
Separation, the collar slopes, and
\eqref{four:maximal:bounds:trace-weighted}--\eqref{four:maximal:bounds:boundary-absorption}
were proved before introducing any arbitrary $C(N)$.
The graph comparisons, Moser iteration, and subsequent regularity
may have arbitrary fixed-$N$ losses. The final mass comparison uses
the earlier weighted trace and boundary-flux inequalities directly;
none of these later constants enters its limit as $N\to\infty$.
\end{remark}

\section{Local regularity of the coupled system}\label{four:maximal:sec:regularity}

The bounded-range estimate makes the trace equation uniformly elliptic
at fixed $N$.  The second equation still contains $|\grad Z|^2$ and
$|\Hess f|^2$, so classical estimates cannot yet be applied to it.
We first use Theorem~\ref{four:thm:axial-regularity} to control $df$
and the local integral of $|\Hess f|^2$.  The scalar measure-source
estimate of Lemma~\ref{b:lem:natural-growth} then gives H\"older
continuity of $Z$, including at both boundary types.  These two gains
start the classical bootstrap.

All constants in this section may depend arbitrarily on the fixed data
and $N$.  The earlier polynomial bounds will control the final energy
error.  In coordinate calculations $D$ denotes ordinary differentiation;
the graph scalar $D=1+l^2|\grad f|^2$ retains its earlier meaning in
geometric formulas.

\begin{proposition}\label{four:maximal:prop:classical-estimates}
Fix the strict data, $\eps$, a sufficiently large $N$, and any of the
homotopy stages of Section~\ref{four:maximal:sec:deformation}.  Every smooth
solution on an allowed finite truncation with $t\ge-\eps$ has uniform
local bounds for every finite number of derivatives of $f$ and $Z$,
including up to the inner and outer boundary faces.  On the uniform
unit patches on the end these constants are independent of the
outer truncation radius and of the stage parameter.  They may depend
arbitrarily on $N$.  In particular, on a fixed finite truncation they
bound the $C^{1,\alpha}$ norm of $Z$ for each fixed $0<\alpha<1$.
\end{proposition}

\begin{proof}
Work on the uniform interior and boundary-normal patches of the
prepared geometry.

\paragraph{The axial estimate and the Hessian measure.}
Proposition~\ref{four:maximal:prop:bounded-range} bounds $f,Z,t,\grad f$ and
keeps $l/\sqrt D$ bounded above and away from zero at fixed $N$.
Dividing the trace equation by this factor gives
\begin{equation}\label{four:maximal:reg:f-normalized}
 \Achi:\Hess f=G_f,\qquad
 G_f=\frac{\sqrt D}{l}
             \bigl(F-\tr_{\Achi}K\bigr),\qquad
 |G_f|\le C(N),\quad \chi\ge\chi _0(N)>0.
\end{equation}
The geometry of each patch has the bounds required by
Theorem~\ref{four:thm:axial-regularity}, and each face value of $f$
is constant.  Each individual solution is smooth; the theorem's
constants depend only on the bounded gradient, source, geometry and
ellipticity.  It gives a uniform H\"older bound for $\grad f$ and
\begin{equation}\label{four:maximal:reg:f-hessian-measure}
 \int_{B_r(x)\cap\Omega_R}|\Hess f|^2\dd V_g
                              \le C(N)r^{2+2\alpha}.
\end{equation}
Coordinate Hessians satisfy the same estimate, since their difference
from covariant Hessians is bounded.  At a face the theorem uses the
$C^1$ odd graph double and controls its interface flux explicitly.
It has no support hypothesis on $K$; the prepared tensor tail and
geometry supply the uniform patch bounds used here.

\paragraph{The scalar source and its boundary reflection.}
In coordinates, absorb the volume density in the second equation and
write it as
\begin{equation}\label{four:maximal:reg:natural-equation}
 \operatorname{div}(\mathcal A DZ+B_1)=E.
\end{equation}
Here $\mathcal A$ and $B_1$ are the coordinate versions of $3u\Achi$
and $\lambda u\Achi K(w,\cdot)$, respectively, with the appropriate
homotopy factors in the source.  On the established range
$\mathcal A$ is bounded, symmetric and uniformly elliptic, and $B_1$
is bounded.  The implicit relation for $t$ has derivative
$1+pv/3>0$, so $t=t(x,Z,Df)$ is smooth on that range and
\begin{equation}\label{four:maximal:reg:t-chain}
 |Dt|\le C(N)(1+|DZ|+|D^2f|).
\end{equation}
The quadratic expression for $\cT$, together with the bounded
lower-order terms in each stage, therefore gives
\begin{equation}\label{four:maximal:reg:coupled-growth}
 |E|\le C(N)(1+|DZ|^2+|D^2f|^2).
\end{equation}
This estimate does not presume a bound for $DZ$.  The prescribed
inner conormal value concerns the total flux and is bounded on this
range.  The outer value of $Z$ is zero.

To apply Lemma~\ref{b:lem:natural-growth} at a boundary, flatten the
face to $x_4=0$, with domain $x_4>0$, and write
$Q=\mathcal A DZ+B_1$.  Put $R=\operatorname{diag}(1,1,1,-1)$.
For $x_4<0$ set
\[
 \widetilde Z(x)=sZ(Rx),\quad
 \widetilde{\mathcal A}(x)=R\mathcal A(Rx)R,\quad
 \widetilde B_1(x)=sRB_1(Rx),\qquad s\in\{1,-1\}.
\]
Then $\widetilde Q(x)=sRQ(Rx)$ below the plane.  Its distributional
divergence has reflected bulk source $sE(Rx)$ and plane source
\begin{equation}\label{four:maximal:reg:reflection-source}
 (1+s)Q_4^+\dd\mathcal H^3\big|_{\{x_4=0\}}.
\end{equation}
At the inner conormal face use $s=1$.  The added signed density is
bounded because $Q_4^+$ is the prescribed total conormal flux in
these coordinates.  At the outer zero Dirichlet face use $s=-1$.
The reflected $Z$ is continuous and the normal total flux matches,
so there is no plane term.  The coefficient matrix remains measurable,
symmetric and uniformly elliptic.  Each reflected solution is continuous
and smooth on its two closed sides.  Integration on those sides gives
the weak exponential product rule required by the shared lemma,
including the full measure in \eqref{four:maximal:reg:reflection-source}.

The absolute source on the doubled patch is consequently dominated by
\[
 C(N)(1+|DZ|^2)\dd x+\mu _0,\qquad
 \mu _0=C(N)|D^2f|^2\dd x
          +2|Q_4^+|\dd\mathcal H^3\big|_{\{x_4=0\}},
\]
with the bulk density reflected and the plane term omitted where absent.
Equation~\eqref{four:maximal:reg:f-hessian-measure} and the area of a
three-dimensional plane section give, at every local center,
\[
 \mu _0(B_r(x))\le C(N)r^{2+\beta},\qquad
                         \beta=\min(2\alpha,1)>0.
\]
On an interior patch only the bulk term is needed.  Thus all the
hypotheses of Lemma~\ref{b:lem:natural-growth}, with dimension four,
hold.  Its H\"older representative agrees with the existing continuous
$Z$, giving a uniform H\"older bound up to either face.

\paragraph{Schauder for the trace equation and the normal datum.}
The actual coefficients in \eqref{four:maximal:reg:f-normalized} are smooth
functions of $x,Z,Df$, including at $Df=0$, since
\[
 \Achi=\Id-\frac{3+p}{3+pv}\,w\otimes w.
\]
Its source is smooth in the same variables and in $f$.  The
H\"older bounds just proved therefore make them uniformly
$C^{0,\theta}$ for some $\theta>0$.  Interior and constant
Dirichlet Schauder estimates give
\begin{equation}\label{four:maximal:reg:f-schauder}
 \|f\|_{C^{2,\theta}(U')}\le C(N)
\end{equation}
on smaller patches; the hypotheses are precisely uniform
ellipticity, $C^{0,\theta}$ leading coefficients and source, smooth
faces and constant Dirichlet data
\cite[Sections~6.1--6.2]{GilbargTrudinger2001}.

Expand the divergence equation for $Z$ using
\eqref{four:maximal:reg:f-schauder}.  A derivative of a coefficient depending on
$x,Z,Df$ is a bounded term plus bounded multiples of $DZ,D^2f$.
Thus it takes the nondivergence form
\begin{equation}\label{four:maximal:reg:Z-nondivergence}
 a^{ij}(x)D_{ij}Z=R_Z,\qquad
 |R_Z|\le C(N)(1+|DZ|^2),
\end{equation}
with uniformly elliptic $C^{0,\theta}$ leading coefficients.  No
derivative of $Z$ of order two has been placed on its right side.

On an inner face, $f$ is constant and $Df$ is normal, so
$\Achi\nu=\chi\nu$.  Dividing the flux condition by $3u\chi>0$
gives
\begin{equation}\label{four:maximal:reg:normal-datum}
 \partial_\nu Z=b(x,Z,f,Df),
\end{equation}
where $b$ is a smooth bounded-range function, separately on each
smooth stage of the homotopy.  The formulas agree at stage junctions
and have uniform bounds.  At zero gradient the smooth formula for
$\Achi$ above gives the same conclusion.  A smooth collar
extension of the right hand side has derivative bounded by
\[
 C(N)(1+|DZ|+|Df|+|D^2f|)
                         \le C(N)(1+|DZ|).
\]
The trace theorem therefore bounds its
$W^{1-1/p,p}$ norm on the face by
$C(N)(1+\|Z\|_{W^{1,p}})$ on a larger collar patch.  In particular,
differentiating this extension costs $DZ$ and $D^2f$, and does not
cost $D^2Z$.

\paragraph{Absorbing the quadratic gradient term.}
Fix $p_1>4$.  The linear local $W^{2,p_1}$ estimates for
\eqref{four:maximal:reg:Z-nondivergence}, with zero Dirichlet data or
\eqref{four:maximal:reg:normal-datum}, give on nested patches
\begin{equation}\label{four:maximal:reg:Z-linear-estimate}
 \|Z\|_{W^{2,p_1}(U')}
 \le C(N)\left(1+\|DZ\|_{L^{2p_1}(U)}^2
                          +\|Z\|_{W^{1,p_1}(U)}\right).
\end{equation}
We specify why the coefficient regularity at this point suffices.
Freeze the continuous leading matrix on small patches.  The constant
coefficient interior Hessian estimate follows from the Laplacian
estimate by a linear coordinate change.  At a boundary, use normal
coordinates.  On that face the matrix has no mixed tangential-normal
entry, by the normal-axis property above.  After removing the
boundary datum by its Sobolev trace extension, odd Dirichlet or even
homogeneous-normal reflection gives the constant coefficient
estimate.  The term
$(a(x)-a(x_0)):D^2Z$ is absorbed when its oscillation is sufficiently
small.  Cutoff terms are first order.  This proves the local
estimate with the $W^{1-1/p_1,p_1}$ normal trace norm just estimated;
it is the scalar Dirichlet/oblique specialization of
\cite[Theorems~15.1--15.3, pp.~702--706, especially
Theorem~15.2, estimate~(15.5)]{AgmonDouglisNirenberg1959}, together with the interior $L^p$ estimate
\cite[Section~9.5]{GilbargTrudinger2001}.  Uniform H\"older control fixes the patch
size uniformly.  The complementing condition can also be checked
directly: a nonzero decaying Fourier mode for the frozen scalar
operator in a half-space has nonzero value and nonzero normal
derivative at its boundary.  Thus neither the Dirichlet nor the
normal operator annihilates that mode.  The normal field is smooth
and has unit transverse component; the boundary is smooth and the
leading coefficients are continuous, as required for this $W^{2,p_1}$
specialization.  The inner and outer faces are disjoint, so no
change of boundary type occurs within such a patch.

For a ball or smooth half-patch $Q_h$ of radius $h$, subtract a
constant from $Z$.  Nirenberg's interpolation inequality, with
derivative orders $1,2$, endpoint exponent $\infty$, Hessian
exponent $p_1$ and interpolation parameter $1/2$, gives
\begin{equation}\label{four:maximal:reg:nirenberg-scaled}
 \|DZ\|_{L^{2p_1}(Q_h)}^2
 \le C\operatorname{osc}_{Q_h}Z\,
                       \|D^2Z\|_{L^{p_1}(Q_h)}
       +Ch^{4/p_1-2}(\operatorname{osc}_{Q_h}Z)^2.
\end{equation}
This is the bounded-domain form of
\cite[Lecture~II, (2.2) and Remark~5,
pp.~125--126]{Nirenberg1959}; extension on half-patches preserves
the estimate.  The power of $h$ follows by rescaling from dimension
four.  Our choice $p_1>4$, $j=1$, $m=2$ avoids the exceptional
endpoint cases in that inequality.

The H\"older estimate for $Z$ makes its oscillation at most $Ch^\theta$.
Cover each larger patch in \eqref{four:maximal:reg:Z-linear-estimate} by such
$Q_h$ with bounded overlap.  Taking the $\ell^{p_1}$ norm of
\eqref{four:maximal:reg:nirenberg-scaled} over this cover gives
\[
 \|DZ\|_{L^{2p_1}(U)}^2
 \le C\sup_{Q_h}\operatorname{osc}_{Q_h}Z\,
                 \|D^2Z\|_{L^{p_1}(U^+)}+C(N,h),
\]
where $U^+$ is a fixed enlargement.  The number of small patches
may affect the harmless constant $C(N,h)$; their overlap, hence the
coefficient of the Hessian norm, is bounded independently of the
truncation.  Choose $h$ so that this coefficient, including the
constant in \eqref{four:maximal:reg:Z-linear-estimate} and the fixed enlargement
factor, is less than $1/4$.  Ordinary interpolation handles the
$W^{1,p_1}$ term with another coefficient less than $1/4$.

To make the absorption on enlarged patches precise, use a uniform
radius cover of the whole finite truncation and let $S$ be the
supremum of the local $W^{2,p_1}$ norms on its smaller patches.
Each fixed enlargement is covered by a bounded number of these
patches, independently of the total number of patches.  The last
two estimates give $S\le C(N)+\tfrac12S$.  For each individual
smooth solution on a finite truncation $S$ is finite before this
inequality is used.  Thus $S\le2C(N)$ uniformly.  The same argument
on nested compact patches gives the local estimate when no global
cover is needed.

\paragraph{Bootstrap and the continuation interface.}
Since $p_1>4$, the bound just proved gives
$Z\in C^{1,\theta'}$ for some $\theta'>0$.  The trace equation now
has $C^{1,\theta''}$ coefficients and source after
\eqref{four:maximal:reg:f-schauder}, so one differentiation and Dirichlet
Schauder give $f\in C^{3,\theta''}$, with
$0<\theta''\le\min(\theta,\theta')$.  Equation
\eqref{four:maximal:reg:Z-nondivergence} consequently has a
$C^{0,\theta''}$ source, and \eqref{four:maximal:reg:normal-datum} has a
$C^{1,\theta''}$ datum.  The scalar oblique/Dirichlet Schauder
estimate gives $Z\in C^{2,\theta''}$; the boundary estimate is
\cite[Theorem~7.3, estimate~(7.8), p.~668]{AgmonDouglisNirenberg1959}, with the
complementing condition verified above.
Repeating these two steps in this order raises the differentiability
by one each time, because the trace equation contains $Z$ without
its derivatives and the second equation contains at most $D^2f$
and $DZ$ in its lower order terms.  The normal datum has the same
derivative count.  Smooth prepared geometry and smooth data complete
the bootstrap to every finite order.

For the scalar solve with trial $Z\in C^{1,\alpha}$ used in
Section~\ref{four:maximal:sec:existence}, this argument makes no assertion that
$D^2Z$ exists.  Its preliminary individual third derivatives for
$f$ require only one differentiation of the trace equation, hence
$DZ$.  After the scalar gradient and axial estimates, bounded
$D^2f$ makes $Df$ Lipschitz; smooth compositions with the trial
$C^{1,\alpha}$ input then give $C^{0,\alpha}$ scalar coefficients,
Dirichlet $C^{2,\alpha}$ estimates, and one more differentiation
gives $C^{3,\alpha}$.  The independent scalar gradient bound and
solvability are proved there.  Finally, for actual coupled solutions
the uniform $C^2$ bound already controls every fixed
$C^{1,\alpha}$ norm, $\alpha<1$, regardless of the first H\"older
exponent obtained above.

The prepared end has uniform smooth geometry on unit patches.
Large outer coordinate spheres admit boundary normal patches with
the same bounds; their curvature and higher local chart derivatives
are bounded uniformly.  All the local constants and covering
overlap numbers in this proof can therefore be chosen independently
of the outer radius.  This proves the proposition.
\end{proof}

\section{Scalar solvability and compact continuation}\label{four:maximal:sec:existence}

Fix the strict data and $0<\eps<1$.  In this section $N\ge4$ and an
allowed finite truncation $\Omega_R$ are fixed.  The boundary components
$B$ and $S_R$ are smooth and disjoint.  Fix also $0<\alpha<1$, with no
requirement that $\alpha$ equal the exponent in
Theorem~\ref{four:thm:axial-regularity}.  We parametrize the four successive
stages of Section~\ref{four:maximal:def:homotopy} by $s\in[0,4]$: $s=0$ is the
desired system and $s=4$ is the final zero-scale system.  This parameter
is unrelated to a collar coordinate.  Constants in the scalar
construction below may depend on $N,\eps,R,\alpha$ and on a specified
bound for the trial input.  They need not be polynomial in $N$.

\subsection{The trace equation on unrestricted trial inputs}

\begin{proposition}[Scalar solution operator]\label{four:maximal:prop:scalar-solve}
For every $Z\in C^{1,\alpha}(\overline\Omega_R)$ and every homotopy
parameter $s$, the trace equation with its assigned constant Dirichlet
values has a unique solution $f=f_s[Z]\in
C^{3,\alpha}(\overline\Omega_R)$.  No lower bound on
$t=t(Z,\grad f)$ is assumed.  On bounded sets of trial inputs, these
solutions have bounded $C^{3,\alpha}$ norms, uniformly in $s$.  The
map $(s,Z)\mapsto f_s[Z]$ is continuous into $C^{3,\alpha}$ and is
locally obtained by inversion of the scalar Dirichlet operator.
\end{proposition}

\begin{proof}
Write the trace prescription at parameter $s$ as
\[
 F_s=\eta_N f+F_s^0(x,Z,\grad f),\qquad
 Q=\frac{\sqrt D}{l},\qquad
 G_s=Q\bigl(F_s-\tr_{\Achi}K\bigr).
\]
Here $F_s^0$ is smooth in its displayed finite-dimensional arguments.
For example, in stage three it is
$(1-\zeta)(C+D_0)+\zeta(\tr_{\Achi}K+D_1)$.
The functions $D_0,D_1$ are evaluated on $w$, whose length is less than
one.  Thus $F_s^0-\tr_{\Achi}K$ is bounded on the finite truncation,
uniformly in the gradient and in $s$.  The normalized scalar equation is
\begin{equation}\label{four:maximal:exist:normalized-trace}
 \Achi:\Hess f=G_s(x,Z,f,\grad f).
\end{equation}
Its matrix is independent of the value of $f$, and
$\partial_fG_s=\eta_N Q>0$.

For the scalar method of continuity introduce a further parameter
$\theta\in[0,1]$, retaining the assigned Dirichlet values and using
\begin{align}
 A^{(\theta)}&=(1-\theta)\Id+\theta\Achi
       =\Id-(1-\chi_\theta)e_f\otimes e_f,
       &\chi_\theta&=1-\theta+\theta\chi,\label{four:maximal:exist:scalar-interpolation}\\
 G_s^{(\theta)}&=(1-\theta)\eta_N f+\theta G_s,
       &A^{(\theta)}:\Hess f&=G_s^{(\theta)}.
       \label{four:maximal:exist:interpolated-equation}
\end{align}
At zero gradient these expressions are interpreted by their smooth
tensor formulas.  In particular,
\begin{equation}\label{four:maximal:exist:strict-height}
 \partial_fG_s^{(\theta)}
       =\eta_N\bigl(1-\theta+\theta Q\bigr)>0.
\end{equation}
Every trace stage has coefficient one on $h=\eta_N f$, so this strict
sign holds throughout both parameters.

\paragraph{Height and large-gradient behavior.}
At a critical point of $f$, $w=0$, $\Achi=\Id$, $t=Z$,
$D_0(x,0)=D_1(x,0)=0$, and $\tr K=0$.
The signs at extrema used in Lemma~\ref{four:maximal:lem:height-end} therefore
persist in \eqref{four:maximal:exist:interpolated-equation}: a fixed interval
containing the assigned boundary values and the zeros of its scalar
right side contains every solution.  Denote its resulting bound by
\begin{equation}\label{four:maximal:exist:scalar-height}
 |f|\le M_f.
\end{equation}
This argument uses neither the divergence equation nor the floor.

Suppose $\|Z\|_{C^0}\le M_Z$ and put $\sigma=|\grad f|$.
The implicit relation $Z=t+\frac16\log(1+l(t)^2\sigma^2)$ implies
$t\to-\infty$ uniformly as $\sigma\to\infty$: if $t$ stayed bounded
below, its right side would tend to infinity.  Eventually $p=N$ and
$l=e^{Nt}$, whence exactly
\begin{equation}\label{four:maximal:exist:large-gradient-identity}
 e^{6Z}=e^{6t}+e^{(6+2N)t}\sigma^2.
\end{equation}
Since the first term tends uniformly to zero and $e^{6Z}$ stays between
two positive constants, this gives
\begin{equation}\label{four:maximal:exist:large-gradient-asymptotics}
 \begin{split}
 t&=-\frac{\log\sigma}{N+3}+O(1),\\
 d=e^{6(t-Z)}&\asymp\sigma^{-6/(N+3)},\qquad
 \chi=\frac{3d}{3+N(1-d)}\asymp d,\\
 Q=\sqrt{\sigma^2+l^{-2}}&\asymp\sigma.
 \end{split}
\end{equation}
The constants here are uniform for bounded $Z$ at the fixed parameters.
On bounded gradient sets the implicit equation has bounded $t$ and
positive $l,d,\chi$.  Since $N\ge4$ makes $6/(N+3)<1$, combining the
two ranges yields constants $c_*>0$ and $C_*>0$ such that
\begin{equation}\label{four:maximal:exist:linear-growth}
 \chi_\theta\ge\chi\ge\frac{c_*}{1+\sigma},\qquad
 |G_s^{(\theta)}|\le C_*(1+\sigma)
       \quad\text{when } |f|\le M_f.
\end{equation}
No uniform ellipticity has yet been asserted.

\paragraph{Logarithmic collar barriers.}
Extend the metric smoothly across all faces to a compact enlargement.
Choose a short distance scale below its injectivity radius, the tubular
radii of the faces, and their positive mutual separation.  The geometry
of all these collars is bounded.  For the inward distance $r$ from a
face with value $b$, consider
\begin{equation}\label{four:maximal:exist:log-barrier}
\begin{gathered}
 b\pm\psi(r),\qquad
 \psi(r)=\frac1k\log(1+B_0r),\\
 q_1(r)=\psi'(r)=\frac{B_0}{k(1+B_0r)},\qquad
 \psi''=-kq_1^2.
\end{gathered}
\end{equation}
Here $k$ and $B_0$ are constants, not any of the geometric tensors.
At a comparison contact the gradient is $\pm q_1\grad r$.
The distance Hessian vanishes in that axial direction, and the three
transverse eigenvalues of $A^{(\theta)}$ are one.  Its contraction on
the positive test is consequently
\[
 \chi_\theta\psi''+
      q_1\sum_{a=1}^3\Hess r(E_a,E_a),
\]
where $E_1,E_2,E_3$ are orthonormal tangent directions to the distance
leaf.  If $q_1\ge1$, \eqref{four:maximal:exist:linear-growth} bounds this above by
$-c_*kq_1/2+Cq_1$; the negative test has the reverse bound.
Choose $k$ so large that this axial term dominates both the distance
Hessian bound and $C_*(1+q_1)$, and also the pair-comparison constants
below.  Then choose a width $\delta_2<1/(4k)$ below all the preceding
geometric scales.  Finally choose $B_0$ large enough that
\begin{equation}\label{four:maximal:exist:barrier-choices}
 q_1(\delta_2)\ge2,\qquad \psi(\delta_2)>2M_f+1.
\end{equation}
These choices are compatible: as $B_0\to\infty$ the two quantities
tend respectively to $1/(k\delta_2)>4$ and to infinity.

At a positive interior maximum of $f-b-\psi(r)$ the solution and test
have the same gradient and the solution Hessian is no larger.  Evaluating
the source at the \emph{solution} height, which obeys
\eqref{four:maximal:exist:scalar-height}, the preceding strict upper bound
contradicts \eqref{four:maximal:exist:interpolated-equation}.
At the outer end of the collar the height inequality follows from
\eqref{four:maximal:exist:barrier-choices}, and at its inner end there is equality.
The negative test is handled with the reversed inequalities.  Thus
\begin{equation}\label{four:maximal:exist:collar-comparison}
 |f(x)-b|\le\psi\bigl(\dist(x,\text{face})\bigr)
       \qquad (\dist(x,\text{face})\le\delta_2).
\end{equation}
The tests themselves need not stay in the bounded height interval;
the source estimate was used only at solution heights.

\paragraph{The full two-point comparison.}
The face barriers control graph slopes only at the boundary.  An interior
two-point comparison supplies the missing global gradient bound.  Its
key feature is that the three transverse eigenvalues agree at both
points; the two axial eigenvalues need only have the same lower bound.

Let $r(x,y)$ be distance in the smooth extension.  On the compact set
$\{(x,y)\in\overline\Omega_R^2:r(x,y)\le\delta_2\}$ consider
\[
 f(x)-f(y)-\psi(r(x,y)).
\]
It is zero on the diagonal and negative when $r=\delta_2$.
It cannot have a positive maximum with a boundary endpoint: if, for
example, $x$ belongs to a face of height $b$, then $y$ lies in that
same face's collar, its distance to the face is at most $r(x,y)$,
and \eqref{four:maximal:exist:collar-comparison} and monotonicity of $\psi$ give
$b-f(y)\le\psi(r(x,y))$.  The case with $y$ on a face is identical.
The separation choice excludes different faces at such distances.

Suppose then that a positive maximum occurs at an interior pair with
$0<r<\delta_2$.  The short geodesic $\gamma$ from $y$ to $x$ is unique
in the extension; the distance is smooth at this pair.  Write
$e_y=\dot\gamma(0)$ and $e_x=\dot\gamma(r)$, with unit speed.  First
variation gives
\begin{equation}\label{four:maximal:exist:pair-gradients}
 \grad f(x)=q_1e_x,\qquad \grad f(y)=q_1e_y,
       \qquad q_1=\psi'(r)\ge1.
\end{equation}
Parallel translate an orthonormal transverse frame $E_1,E_2,E_3$
along $\gamma$.  Vary both endpoints simultaneously in the corresponding
$E_a$ directions, using geodesic endpoint curves.  The first variation
of the distance is zero.  The second variation is bounded above by
$Cr$: the parallel comparison field has index form
\[
 -\int_0^r
 \bigl\langle\operatorname{Rm}(E_a,\dot\gamma)\dot\gamma,E_a
 \bigr\rangle\dd v,
\]
whose absolute value is at most $Cr$, and the Jacobi field with the
same endpoints minimizes this index form on the short segment.
The second derivative test at the pair maximum therefore gives
\begin{equation}\label{four:maximal:exist:pair-transverse}
 \sum_{a=1}^3\bigl[
  \Hess f_x(E_a(r),E_a(r))-
  \Hess f_y(E_a(0),E_a(0))\bigr]\le Cq_1r.
\end{equation}
Varying only $x$ in direction $e_x$, and only $y$ in direction $e_y$,
the distance changes linearly and its second derivative vanishes.
Hence
\[
 \Hess f_x(e_x,e_x)\le\psi'',\qquad
 \Hess f_y(e_y,e_y)\ge-\psi''.
\]
The two axial coefficients need not be equal.  Positivity and
\eqref{four:maximal:exist:linear-growth} nevertheless imply
\begin{equation}\label{four:maximal:exist:pair-axial}
 \begin{split}
 &\chi_\theta(x)\Hess f_x(e_x,e_x)
        -\chi_\theta(y)\Hess f_y(e_y,e_y)\\
 &\hspace{12mm}\le
    \bigl(\chi_\theta(x)+\chi_\theta(y)\bigr)\psi''
    \le-\frac{2c_*kq_1^2}{1+q_1}\le-c_*kq_1.
 \end{split}
\end{equation}
There are exactly three transverse eigenvalues at \emph{each} endpoint,
all equal to one.  Thus \eqref{four:maximal:exist:pair-transverse} and
\eqref{four:maximal:exist:pair-axial}, followed by the two scalar equations, give
\[
 G_s^{(\theta)}(x)-G_s^{(\theta)}(y)
       \le Cq_1r-c_*kq_1.
\]
On the other hand, the source growth bound and
\eqref{four:maximal:exist:pair-gradients} give a lower bound
$-2C_*(1+q_1)\ge-4C_*q_1$.  Taking $k$ above so that
$c_*k>C+4C_*$ and $\delta_2<1$ is a contradiction.
This uses no continuity estimate on the two axial coefficients or on
their difference.  The connecting segment is permitted to leave
$\Omega_R$: only endpoint variations and the extension's bounded
geometry enter the argument, and the endpoints are interior.

It follows that $|f(x)-f(y)|\le\psi(r(x,y))$ for short pairs.  Letting
$y\to x$, and also using the collar comparison at a face, proves
\begin{equation}\label{four:maximal:exist:scalar-gradient}
 \|\grad f\|_{C^0(\overline\Omega_R)}\le\psi'(0)=B_0/k.
\end{equation}
This estimate is uniform in $s,\theta$ and on bounded trial-input sets,
and was obtained without assuming the floor or a positive uniform
ellipticity constant in advance.

\paragraph{Regularity and the scalar method of continuity.}
The gradient bound and the implicit relation now confine all scalar
coefficients to a compact range.  In particular
$0<\chi_0\le\chi_\theta\le1$, and the normalized source is bounded.
There is a small regularity point before applying
Theorem~\ref{four:thm:axial-regularity}.  An individual
$C^{2,\alpha}$ solution of \eqref{four:maximal:exist:interpolated-equation} is
$C^{3,\alpha}$ for a trial $Z\in C^{1,\alpha}$.
Indeed one differentiation of the equation gives a linear equation
for a first derivative of $f$.  Its coefficients and source involve
$Z,\grad Z,\grad f,\Hess f$ and smooth geometric coefficients;
they involve no second derivative of $Z$.
Difference quotients justify this differentiation initially.
In a flattened boundary chart tangential differentiation preserves
the zero Dirichlet value after subtracting the constant face height.
The linear boundary Schauder estimate controls these tangential
derivatives through order two.  Solving the original equation for
the pure normal second derivative, whose coefficient is positive,
and differentiating that identity in the normal direction recovers
the remaining pure normal third derivative.  Every other third
derivative already has a tangential index.  This gives individual
third-order regularity, without presuming a uniform estimate.

Theorem~\ref{four:thm:axial-regularity} now applies to
\eqref{four:maximal:exist:interpolated-equation}, including its constant-Dirichlet
faces.  It bounds $\grad f$ in $C^{0,\beta}$ for some $\beta>0$.
Since the trial inputs are bounded in $C^{1,\alpha}$, the coefficients
and source of the scalar equation are consequently bounded in
$C^{0,\beta'}$ for a positive $\beta'$.  Dirichlet Schauder estimates
give a uniform $C^{2,\beta'}$ bound for $f$.
In particular $\grad f$ is now uniformly Lipschitz.  Smooth composition
with $Z$, which also has bounded first derivatives, shows that those
same coefficients and sources have bounded $C^{0,\alpha}$ norms for
the originally chosen, arbitrary $\alpha<1$.
Schauder estimates upgrade $f$ to $C^{2,\alpha}$.
Their coefficients and sources then have bounded $C^{1,\alpha}$
norms, and the preceding differentiation argument gives
\begin{equation}\label{four:maximal:exist:scalar-schauder}
 \|f\|_{C^{3,\alpha}(\overline\Omega_R)}\le C.
\end{equation}
The linear estimates used here are the interior and Dirichlet boundary
Schauder estimates and their differentiated versions
\cite[Sections~6.1--6.4]{GilbargTrudinger2001}.  Their hypotheses hold on the finite
smooth domain: positive uniform ellipticity on the established range,
the indicated H\"older coefficient bounds, and smooth constant
Dirichlet data on each disjoint component.

For completeness, the local inverse has the required sign.
Linearizing $A^{(\theta)}:\Hess f-G_s^{(\theta)}$ in the $f$
variable gives
\begin{equation}\label{four:maximal:exist:scalar-linearization}
 \mathcal L\varphi
    = A^{(\theta)ij}\nabla_i\nabla_j\varphi
           +b^i\nabla_i\varphi
           -\eta_N(1-\theta+\theta Q)\varphi,
      \qquad \varphi|_{\partial\Omega_R}=0.
\end{equation}
Its first-order coefficients are bounded H\"older functions along each
solution.  Differentiating the gradient dependence of $A^{(\theta)}$
contributes terms involving $\Hess f$ only to this first-order part.
The strictly negative zeroth-order coefficient and the maximum
principle give a zero kernel.  The linear Dirichlet solvability theorem
and Schauder estimate give an isomorphism
$C^{2,\alpha}_0\to C^{0,\alpha}$
\cite[Sections~6.2--6.3]{GilbargTrudinger2001}; here the subscript means zero
trace on \emph{all} boundary components.  Equivalently, one can
continue the linear operator to $\Delta-\eta_N$, using its Schauder
estimate and the maximum principle at each point of that linear
continuation.  With the additional coefficient regularity it is also
an isomorphism $C^{3,\alpha}_0\to C^{1,\alpha}$.

The endpoint $\theta=0$ is the solvable linear Dirichlet problem
$\Delta f=\eta_N f$.  Subtract a smooth extension of the assigned face
values when working on spaces with homogeneous boundary data.
The implicit function theorem and \eqref{four:maximal:exist:scalar-linearization}
give openness in $\theta$.  The bound
\eqref{four:maximal:exist:scalar-schauder} gives compactness in $C^{2,\alpha}$,
so passage to a limiting equation gives closedness.  Thus the solution
exists at $\theta=1$.

If $f_1-f_2$ had a positive maximum in the interior, their gradients
and therefore their principal matrices and $Q$ would agree there.
Their Hessian difference would have nonpositive contraction, whereas
the difference of the sources would be
$\eta_N Q(f_1-f_2)>0$.  This contradiction, and its reverse, prove
uniqueness.  Finally use \eqref{four:maximal:exist:scalar-linearization} in
$C^{3,\alpha}_0\to C^{1,\alpha}$ with $Z$ as a
$C^{1,\alpha}$ parameter.  Smooth composition and the implicit
function theorem show that the unique scalar solution depends
continuously on that input in $C^{3,\alpha}$.  The same statement
holds along each homotopy stage; the explicit prescriptions and
boundary heights agree at their junctions, so it holds across those
junctions as well.
\end{proof}

\subsection{The frozen linear return problem}

The trace equation is now solvable for every trial $Z$, including trials
below the intended floor.  We can therefore define a compact return map
on a fixed function space.  The floor becomes a restriction on the
degree domain, rather than an unproved prerequisite for defining the map.

Set
\[
 X=\{Z\in C^{1,\alpha}(\overline\Omega_R):Z|_{S_R}=0\}.
\]
Given $(s,Z)\in[0,4]\times X$, first compute $f=f_s[Z]$ by
Proposition~\ref{four:maximal:prop:scalar-solve} and then all of $t,l,u,w,\chi$,
$F$ and $\cT$ from this pair.  Let $\lambda_s$ denote the drift
parameter, let $q_s=1$ in the first three stages and $q_s=\eta$ in
the last, and let $b_s$ denote the prescribed right side of
$(V_{\lambda_s})_\nu/u$ in the full homotopy, including its last-stage
scale.  Define $T_sZ=\widetilde Z$ by
\begin{equation}\label{four:maximal:exist:return-problem}
 \begin{cases}
 \divg\bigl(3u\Achi\grad\widetilde Z
                  +u\Achi\lambda_sK(w,\cdot)\bigr)=q_su\Xi
                         &\text{in }\Omega_R,\\
 \bigl(3u\Achi\grad\widetilde Z
                  +u\Achi\lambda_sK(w,\cdot)\bigr)_\nu=ub_s
                         &\text{on }B,\\
 \widetilde Z=0          &\text{on }S_R.
 \end{cases}
\end{equation}
Every quantity in this display except $\grad\widetilde Z$ is frozen
at the input pair, including every occurrence of $\grad Z$ in $\Xi$.

The coefficient $\mathcal A=3u\Achi$ is symmetric and positive
definite.  On a bounded subset of $X$ its ellipticity constants are
uniformly positive, by Proposition~\ref{four:maximal:prop:scalar-solve} and the
smooth implicit equation for $t$.  Let
$\mathcal D=u\Achi\lambda_sK(w,\cdot)$.
On the Hilbert space of $H^1$ functions vanishing on $S_R$, the weak
form of \eqref{four:maximal:exist:return-problem} is
\begin{equation}\label{four:maximal:exist:return-weak}
 \int_{\Omega_R}\mathcal A\grad\widetilde Z\cdot\grad\varphi
 =-\int_{\Omega_R}q_su\Xi\varphi
  -\int_{\Omega_R}\mathcal D\cdot\grad\varphi
  -\int_Bub_s\varphi.
\end{equation}
The last sign uses the inward normal $\nu$ on $B$.
Poincar\'e's inequality holds because $S_R$ is nonempty and the
domain is connected.  Consequently the left side is a coercive
bilinear form and the right side is a bounded linear functional,
by the trace inequality.  To solve it directly, minimize one half
of this quadratic form minus the right-side functional on that
Hilbert space.  Coercivity bounds a minimizing sequence in $H^1$;
weak compactness and lower semicontinuity give a minimizer.  Its first
variation is \eqref{four:maximal:exist:return-weak}, and strict convexity gives
uniqueness.  There is no Neumann compatibility constraint.

We check the precise regularity of this linear problem.  The formula
\[
 \Achi=\Id-\frac{3+p}{3+pv}\,w\otimes w
\]
is smooth also at zero gradient, and the derivative of the implicit
relation in $t$ is $1+pv/3>0$.
Thus $t,u,\Achi,\mathcal D$ have bounded $C^{1,\alpha}$ norms on
bounded input sets.  Differentiating $t=t(x,Z,\grad f)$ requires only
$\grad Z$ and $\Hess f$.  The expression for $\cT$ is quadratic in
$\Hess f$ and $\grad t$ with smooth bounded-range coefficients,
so $q_su\Xi$ has a bounded $C^{0,\alpha}$ norm.  The prescribed
boundary expression $ub_s$ has a bounded $C^{1,\alpha}$ norm.
At $B$ the scalar Dirichlet data make $\grad f$ normal, hence
$\Achi\nu=\chi\nu$; the boundary equation is equivalently
\begin{equation}\label{four:maximal:exist:return-normal}
 \partial_\nu\widetilde Z
  =\frac{b_s-\lambda_s(\Achi K(w,\cdot))_\nu}{3\chi}.
\end{equation}
This is a $C^{1,\alpha}$ normal derivative datum, with a denominator
bounded away from zero on the input set.
After expanding the divergence, interior and boundary Schauder
estimates therefore give a bounded $C^{2,\alpha}$ norm for
$\widetilde Z$.  The estimates apply locally with Dirichlet data on
$S_R$ and a uniformly oblique, here normal, derivative on $B$;
these faces do not meet.  We use the Dirichlet estimates just cited
and the oblique derivative estimates of
\cite[Theorem~6.30, p.~127]{GilbargTrudinger2001}; the weak existence argument above fixes
the additive constant that is present in a pure Neumann problem.

Subtracting two return equations and applying the same linear
estimates proves continuity in $(s,Z)$ into $C^{2,\alpha}$.
The parameter continuity is uniform on bounded input sets.  Indeed
the negative zeroth-order bound in \eqref{four:maximal:exist:scalar-linearization}
and the Schauder estimate bound the scalar inverse uniformly on
such a set.  Differentiation in a stage parameter, after extending
its Dirichlet data, then bounds the parameter derivative of $f_s[Z]$;
the return equation gives the same conclusion for $T_sZ$.
The finitely many stage junctions preserve continuity.
The inclusion $C^{2,\alpha}\hookrightarrow C^{1,\alpha}$ on the
fixed smooth compact domain is compact.  We have thus constructed
a continuous homotopy $T_s:X\to X$ that is compact on bounded sets,
jointly in $s$, before making any floor restriction.
A fixed point satisfies both original equations and their exact
boundary conditions.  It is smooth: the first gain is
$Z\in C^{2,\alpha}$, and successive differentiation of the scalar
equation, the divergence equation and its normal boundary condition
then alternately improves $f$ and $Z$ to every finite order.

\subsection{Degree on the open floor sections}

\begin{proposition}[Existence and exhaustion]\label{four:maximal:prop:existence}
For the fixed strict data and every $0<\eps<1$ there are $N_0$ and
$R_{\min}(N)$, with $R_{\min}(N)\to\infty$, such that for every
integer $N\ge N_0$ and every $R\ge R_{\min}(N)$ the desired coupled
system has a smooth solution on $\overline\Omega_R$ with
$t>-\eps$.  At each such fixed $N$, a subsequence as $R\to\infty$
converges smoothly on compact subsets of $\overline\Omega$ to a
solution of the equations and inner boundary conditions with
$t\ge-\eps$.  The previously proved height, range and collar
estimates hold for this limit.
\end{proposition}

\begin{proof}
All estimates from Sections~\ref{four:maximal:sec:deformation}--\ref{four:maximal:sec:regularity}
were conditional estimates for arbitrary smooth homotopy solutions.
Specifically, Lemma~\ref{four:maximal:lem:height-end} and Proposition~\ref{four:maximal:prop:floor}
precede the bounded-range and derivative estimates;
Lemma~\ref{four:maximal:lem:separation} and Lemma~\ref{four:maximal:lem:collar-slopes} apply to
solutions satisfying the floor;
Proposition~\ref{four:maximal:prop:bounded-range} then bounds their variables;
and Proposition~\ref{four:maximal:prop:classical-estimates} bounds their
derivatives.  None presupposes a fixed point of the map just defined.
On this finite truncation these results give, uniformly in $s$,
\[
 \|Z\|_{C^{1,\alpha}}\le M_*
 \quad\text{for every smooth fixed point with }\min t\ge-\eps.
\]
The arbitrary exponent $\alpha<1$ is permitted because the classical
estimates give a $C^2$ bound, and higher bounds as needed.  Choose
$M_0>M_*+1$.

Let
\begin{equation}\label{four:maximal:exist:floor-sections}
 \begin{split}
 m(s,Z)&=\min_{\overline\Omega_R}t\bigl(x,Z,\grad f_s[Z]\bigr),\\
 \mathcal W&=\{(s,Z)\in[0,4]\times X:
            m(s,Z)>-\eps, \|Z\|_{C^{1,\alpha}}<M_0\},\\
 \mathcal W_s&=\{Z:(s,Z)\in\mathcal W\}.
 \end{split}
\end{equation}
The minimum is continuous by Proposition~\ref{four:maximal:prop:scalar-solve},
so $\mathcal W$ is relatively open and every section is a bounded
open subset of $X$.  Consider
\[
 \mathcal K=
 \{(s,Z)\in\overline{\mathcal W}:Z=T_sZ\}.
\]
This is compact: any sequence has a convergent parameter subsequence,
and joint compactness of $T_s$ on the bounded input ball makes its
fixed points converge in $X$; continuity preserves the fixed-point
equation and membership in the closed set.
Every point of $\mathcal K$ is a smooth solution with
$m(s,Z)\ge-\eps$.  Its norm is at most $M_*<M_0$, and
Proposition~\ref{four:maximal:prop:floor} excludes $m(s,Z)=-\eps$.
Therefore
\begin{equation}\label{four:maximal:exist:no-boundary-fixed-point}
 \mathcal K\subset\mathcal W.
\end{equation}

Apply Lemma~\ref{b:lem:moving-degree} with $I=[0,4]$, the Banach
space $X$ above, the jointly compact homotopy $(s,Z)\mapsto T_sZ$, and
$\mathcal U=\mathcal W$.  The uniform norm bound makes $\mathcal W$
bounded; continuity of $m$ makes it relatively open.  The preceding
compactness argument and \eqref{four:maximal:exist:no-boundary-fixed-point}
verify that its closed fixed-point set is compact and contained in
$\mathcal W$, so no fixed point lies on the relative boundary.
The lemma gives constancy of
\[
 \deg(\Id-T_s,\mathcal W_s,0)\qquad (0\le s\le4).
\]
It applies to these varying sections without a convexity assumption.

At $s=4$ the scalar boundary values are zero and the final trace
equation is
\[
 \tr_{\Achi}H^f=\eta_N f+D_1(x,w).
\]
At a positive or negative extremum $w=0$ and $D_1(x,0)=0$, so
the strict height sign forces $f=0$.  This holds for every trial
$Z$, not just at fixed points.  Thus $t=Z$, $\chi=1$ and $u=L_0(Z)$.
Both the drift and the prescribed source and boundary flux in
\eqref{four:maximal:exist:return-problem} vanish.  Testing its homogeneous
equation by $\widetilde Z$ gives
\[
 \int_{\Omega_R}3L_0(Z)|\grad\widetilde Z|^2=0,
\]
so $T_4Z=0$ for every input.  At its sole fixed point $Z=0$ one
has $t=0$, so $0\in\mathcal W_4$ and
$\deg(\Id-T_4,\mathcal W_4,0)=1$
\cite[Chapter~II, \S8, p.~55]{LeraySchauder1934}.
The same degree at $s=0$ yields a desired fixed point in
$\mathcal W_0$, proving the assertion on the finite truncation.

We finally make the parameter order explicit.  Fix the strict
background, $\eps$, the collars and the end weights first.
The pointwise floor argument chooses the small $\delta_0$ and a
polynomially large penalty coefficient $C_N$ without derivative
estimates.  The separation statement is uniform along arbitrary
expanding sequences of the first two homotopy stages; applied at
the finitely many fixed collar ends, it supplies the eventual
separation needed for the slopes.  Increase $N_0$ to meet this
requirement and the finitely many conditions of the form
$\Pi_N\eta_N/\ell<\text{constant}$.  Increase $R_{\min}(N)$ to
meet all the fixed-$N$ outer-radius requirements, requiring also
$R_{\min}(N)\ge N$.  These choices give the estimates for every
$N\ge N_0$ and every $R\ge R_{\min}(N)$.
All subsequent scalar, degree-radius and classical regularity
constants may depend arbitrarily on the now fixed $N$.

For that fixed $N$, take $R_j\to\infty$ and the desired fixed
points just obtained.  Proposition~\ref{four:maximal:prop:classical-estimates}
gives uniform bounds of every finite order on each fixed compact
subset, including patches meeting $B$.
Successive compact extractions on a countable exhaustion produce
a subsequence converging in $C^k$ on each such compact set for
every finite $k$.  All equations, inner data and collar inequalities
pass to the limit.  The strict floor may become non-strict, which
is the asserted $t\ge-\eps$.
The normalization at infinity is supplied by the uniform tail
estimates in Proposition~\ref{four:maximal:prop:end-control}; it is not inferred
from compact convergence alone.  Only after this fixed-$N$
exhaustion will $N$ tend to infinity in the flux comparison.
The latter uses the earlier polynomial estimates directly, and
does not use the arbitrary constants of the present section.
\end{proof}

\section{End normalization and the mass--area comparison}
\label{four:maximal:sec:comparison}

Fix the strict approximation data, the threshold exterior $\Omega$ with
boundary $B$, and $0<\eps<1$.  Take $N$ sufficiently large as in
Proposition~\ref{four:maximal:prop:existence}.  In the first part of this section, $N$
remains fixed while the truncation radius tends to infinity.  Constants
in end estimates may depend arbitrarily on these fixed parameters.  In
the flux deficit estimate, by contrast, we will return directly to the
polynomial bounds of Lemma~\ref{four:maximal:lem:weighted-trace}.

\subsection{Normalization inherited from the truncations}

\begin{proposition}[Control of the asymptotic end]\label{four:maximal:prop:end-control}
The exhaustion solutions of Proposition~\ref{four:maximal:prop:existence} admit a
smooth local limit $(f,Z)$ on $\overline\Omega$ satisfying the desired
system and $t\ge-\eps$.  Every such limit has
\begin{equation}\label{four:maximal:end:decay}
 |\nabla^j f|\le C_j(N)e^{-c_j(N)r}\quad(j\ge0),
 \qquad Z,t=O_2(r^{-2}).
\end{equation}
The zeroth-order estimates producing this normalization are uniform on
the finite truncations.  The resulting metric
$\ghat=e^{2t}(g+l^2df\otimes df)$ is smooth, complete with its boundary
included, and asymptotically flat with
\begin{equation}\label{four:maximal:end:geometry}
 \ghat_{ij}-\delta_{ij}=O_2(r^{-q}),\qquad
 0\le R_{\ghat}=O(r^{-4-\delta}),\qquad \widehat H_B<0.
\end{equation}
Here the normal defining $\widehat H_B$ points toward infinity.
\end{proposition}

\begin{proof}
The local compactness and the inner boundary conditions follow from
Propositions~\ref{four:maximal:prop:classical-estimates} and
\ref{four:maximal:prop:existence}.  To determine the value at infinity,
apply Lemma~\ref{b:transition:end} to the finite truncations before
taking their local limit.  We verify its hypotheses explicitly.

The strict background has $1<q<2$, $0<\delta<q-1$,
$g-\delta_{\mathrm{Eucl}}=O_2(r^{-q})$,
$K=O_1(r^{-1-q})$, $\tr_gK=0$, and
$R_g=O(r^{-4-\delta})$ by
Proposition~\ref{four:maximal:prop:strict-data}.  Outside the compact
support of $C$, the trace right side is $\eta_N f$.
The divergence equation is exactly
\eqref{b:transition:divergence}, with
$\rho_0=r^{-4-\delta}$, $\rho_1=r^{-2\gamma}$ and the smooth penalty
switch defined in \eqref{four:maximal:def:penalty}.  Their differentiated
end bounds follow from the fixed smooth radial choices.
Propositions~\ref{four:maximal:prop:bounded-range} and
\ref{four:maximal:prop:classical-estimates} supply the common bounded
range and all fixed-$N$ classical estimates, on end unit patches and on
the zero-Dirichlet outer faces, uniformly in the truncation radius.
The leading matrix is uniformly elliptic on that range.

The lemma therefore gives exponential decay of $f$ and its fixed
derivatives, a truncation-uniform bound $|Z|+|t|\le C(N)r^{-2}$, and
$Z,t=O_2(r^{-2})$ for every local limit.  In particular normalization
at infinity follows from a comparison on the actual truncations, not
from local convergence.  The equation used for that comparison is
\begin{align}\label{four:maximal:end:zero-graph}
 3\Delta_g Z
 &+3\bigl[p(Z)+2-\delta_0(p(Z)+1)\bigr]|dZ|^2\notag\\
 &=\tfrac{\delta_0}{2}|K|^2+\rho
       -m_0(Z)C_N\rho_0+\mathcal E_N,
\end{align}
where the graph error $\mathcal E_N$ has exponential unit-patch bounds.
The $|K|^2/2$ term survives because the tensor is retained; its decay is
integrable since $2+2q>4+\delta$.  All constants in this application
are allowed to depend on the fixed $N$.  They will not be used to
estimate the later limit in $N$.

The signs can now be recorded directly from the desired system.
Since $F=\eta_N f+C$ and $w(f)=a\sigma$, the scalar identity becomes
\begin{align}\label{four:maximal:end:scalar-sign}
 \tfrac12e^{2t}R_{\ghat}
 ={}&[\mu+J(w)+w(C)-\rho]
       +(1-\delta_0)\cT\notag\\
    &+(1-\delta_0)\eta_N a\sigma+m_0(t)\mathcal P.
\end{align}
The bracket is nonnegative by $|w|\le1$ and
$\mu-|J|\ge|dC|+\rho$; all remaining terms are nonnegative.
The boundary identity and the prescribed conormal datum give
\begin{equation}\label{four:maximal:end:boundary-sign}
 H_B+3\partial_\nu t=-N_B,
 \qquad \widehat H_B=-e^{-t}\sqrt d\,N_B<0.
\end{equation}
This uses the normal into $\Omega$ on every component of $B$.

As quadratic forms $\ghat\ge e^{-2\eps}g$.  A $\ghat$-Cauchy sequence
is therefore $g$-Cauchy; completeness of the background exterior and
local equivalence of the smooth metrics give completeness with the
boundary included.  Since $q<2$, the conformal $O_2(r^{-2})$ term and
the exponential graph error retain the background decay
$\ghat-\delta=O_2(r^{-q})$.  Equation~\eqref{four:maximal:end:zero-graph} gives
$\Delta_g t=O_N(r^{-4-\delta})$, since $|dt|^2=O_N(r^{-6})$ and
$0<\delta<1$.  Finally,
\[
 R_{\ghat}
  =e^{-2t}\bigl(R_{\gbar}-6\Delta_{\gbar}t
                                  -6|dt|_{\gbar}^2\bigr)
\]
and the exponential graph errors, together with
$R_g=O(r^{-4-\delta})$, give the scalar decay in
\eqref{four:maximal:end:geometry}.
\end{proof}

\subsection{The ADM change and a polynomial flux deficit}

All integrals in the next Proposition use the background metric $g$.
At the inner boundary $\nu$ points into $\Omega$, and at a coordinate
sphere $\nu_g$ points toward increasing radius.

\begin{proposition}[Flux and mass comparison]\label{four:maximal:prop:flux}
For the limiting solution at fixed $N$, the flux
\[
 \mathfrak F=\lim_{r\to\infty}\int_{S_r}V\cdot\nu_g\dd A_g
\]
exists and satisfies
\begin{equation}\label{four:maximal:end:mass-flux}
 E_{\ghat}=E_g-\frac{\mathfrak F}{3\omega_3},
 \qquad \omega_3=2\pi^2.
\end{equation}
There is a polynomial bound $\Pi_N$, independent of the exhaustion
radius and of the fixed-$N$ regularity constants, such that
\begin{equation}\label{four:maximal:end:flux-deficit}
 \mathfrak F\ge-\Pi_N\left(\ell+\frac{\ell^2}{\eta_N}\right).
\end{equation}
Consequently, for fixed strict data and $\eps>0$,
$E_{\ghat}\le E_g+o(1)$ as $N\to\infty$.
\end{proposition}

\begin{proof}
The end estimates imply
\begin{align*}
 \cT&=\tfrac12|K|^2+3(p(t)+1)|dt|^2+O_N(e^{-cr}),\\
 u&=1+O_N(r^{-2}),\qquad
 V=3\grad_g t+O_N(r^{-5})+O_N(e^{-cr}).
\end{align*}
In particular $u\Xi$ is absolutely integrable.  Its nonexponential
terms have decay $r^{-2-2q}$, $r^{-6}$, or $r^{-4-\delta}$.
The possibly slower weight $\rho_1$ is multiplied by the exponentially
small $v$ on this fixed-$N$ end.  The divergence theorem gives
\[
 \int_{S_r}V\cdot\nu_g\dd A_g
    -\int_B V\cdot\nu\dd A_g
       =\int_{\Omega_r}u\Xi\dd V_g.
\]
Here $\Omega_r$ denotes the part of $\Omega$ inside $S_r$.
Absolute integrability proves existence of the limit and yields the
orientation-sensitive formula
\begin{equation}\label{four:maximal:end:divergence-flux}
 \mathfrak F=\int_\Omega u\Xi\dd V_g
                        +\int_B V\cdot\nu\dd A_g.
\end{equation}

We next compute the mass change in the same asymptotic chart.  The
linear conformal perturbation is $2t\delta_{ij}$, and in dimension
four its contribution to the ADM numerator is exactly
\[
 \partial_j(2t\delta_{ij})-\partial_i(2t\delta_{jj})
       =-6\partial_i t.
\]
The differentiated remainders are
$O_N(r^{-3-q})+O_N(r^{-5})+O_N(e^{-cr})$ and have vanishing sphere
integrals.  They include products of $t$ with $g-\delta$, quadratic
conformal terms, and the graph perturbation.  Replacing Euclidean
surface flux for $dt$ by $g$-surface flux changes it by
$O_N(r^{-q})$ after integration.  Since the ADM denominator is
$6\omega_3$, we obtain
\[
 E_{\ghat}-E_g
   =-\frac1{\omega_3}\lim_{r\to\infty}
                  \int_{S_r}\partial_{\nu_g}t\dd A_g
   =-\frac{\mathfrak F}{3\omega_3}.
\]
This also proves existence of the ADM energy of $\ghat$ from the
assumed existence of $E_g$; no interchange of the $r$ and $N$ limits
is involved.

To estimate \eqref{four:maximal:end:divergence-flux}, use the boundary flux bound
\eqref{four:maximal:bounds:boundary-flux} and the weighted trace estimate
\eqref{four:maximal:bounds:trace-weighted} with test function one.  They give
\begin{align}\label{four:maximal:end:boundary-cost}
 \int_B|V\cdot\nu|\dd A_g
 &\le C\frac{\eta_N}{\ell}\int_B L_0\dd A_g\notag\\
 &\le\Pi_N\frac{\eta_N}{\ell}
      \int_{\mathrm{collars}}u(1+\sqrt{\cT})\dd V_g\notag\\
 &\le\Pi_N\frac{\eta_N}{\ell}
      \int_{\mathrm{collars}}u(\rho+\delta_0\cT)\dd V_g.
\end{align}
For the last step, $\rho$ has a positive minimum on these fixed compact
collars, and $\sqrt{\cT}\le c\cT+C_c$ with fixed $c>0$.
Only the polynomial constants in those two earlier estimates enter
\eqref{four:maximal:end:boundary-cost}.  Since $\eta_N/\ell=e^{-\eps N/2}$, increasing
$N$ absorbs this entire cost into, for example, one half of the
$u(\rho+\delta_0\cT)$ contribution to the bulk integral.

The negative bulk term is $um_0(t)\mathcal P$.  On its support,
\[
 -\eps\le t\le-\eps+1/N,
 \qquad \ell\le l\le e\ell,\qquad c\ell\le L_0\le C\ell.
\]
These constants are uniform for $0<\eps<1$ and sufficiently large
$N$.  From the definitions, rather than from a derivative estimate,
one obtains
\begin{equation}\label{four:maximal:end:band-identities}
 u\le C(\ell+\ell^2\sigma),\qquad
 uv\le C\ell^2\sigma,\qquad
 ua\sigma=L_0l\sigma^2\ge c\ell^2\sigma^2.
\end{equation}
For the middle inequality use
$v\sqrt D=l^2\sigma^2/\sqrt{1+l^2\sigma^2}\le l\sigma$.
Thus Young's inequality, with $C_N$ polynomial and $0\le m_0\le1$,
gives
\begin{align}\label{four:maximal:end:penalty-loss}
 um_0\mathcal P
 &\le \Pi_N\bigl[\ell\rho_0+
                         \ell^2\sigma(\rho_0+\rho_1)\bigr]\notag\\
 &\le\tfrac12\delta_0\eta_N ua\sigma
       +\Pi_N\left[\ell\rho_0+
                  \frac{\ell^2}{\eta_N}(\rho_0^2+\rho_1^2)\right].
\end{align}
Off the band the left side vanishes, so the final inequality holds
everywhere.  All weights in the loss are integrable:
\[
 \rho_0=r^{-4-\delta}\in L^1,\qquad
 \rho_0^2\in L^1,\qquad
 \rho_1^2=r^{-4\gamma}\in L^1
\]
on the four-dimensional end, using $\delta>0$ and $\gamma>1$.
Their smooth extensions are bounded on the compact interior.  In
particular no integrability of $\rho_1$ itself is being used.

Insert \eqref{four:maximal:end:boundary-cost} and \eqref{four:maximal:end:penalty-loss} into
\eqref{four:maximal:end:divergence-flux}, absorb the indicated nonnegative terms,
and discard the remaining nonnegative bulk terms.  This proves
\eqref{four:maximal:end:flux-deficit}.  Finally,
\[
 \ell+\frac{\ell^2}{\eta_N}
       =e^{-\eps N}+e^{-\eps N/2},
\]
which tends to zero against every polynomial in $N$.  The constants
used to prove \eqref{four:maximal:end:decay} have not entered this estimate.
\end{proof}

\begin{proof}[Completion of Proposition~\ref{four:maximal:prop:geometric-output}]
The threshold exterior and its all-cut comparison are supplied by
Lemma~\ref{four:maximal:lem:threshold-collars}.  Propositions~
\ref{four:maximal:prop:existence} and \ref{four:maximal:prop:end-control}
construct the complete metric with the stated signs and decay.
Its lower metric bound gives the asserted three-dimensional area
comparison on each cut separately.  Proposition~\ref{four:maximal:prop:flux}
gives the energy estimate, since $\ell=e^{-\eps N}$ and
$\eta_N=\ell^{3/2}$.  Every stated output is therefore established
before the Riemannian Penrose inequality is applied.
\end{proof}

\subsection{A minimal enclosing hypersurface}

The boundary $B$ has strictly negative mean curvature in $\ghat$.
We obtain the minimal boundary required for the Riemannian theorem
from the full-perimeter enclosure proposition.  Perimeter here means the full
area of the measure-theoretic boundary of a set, including any portion
coinciding with an obstacle.

\begin{lemma}[Minimal enclosure]\label{four:maximal:lem:minimal-enclosure}
The metric $\ghat$ of Proposition~\ref{four:maximal:prop:end-control} admits a
nonempty compact smooth embedded minimal hypersurface $\Sigma$
enclosing $B$.  Its exterior is connected, complete with boundary,
and has the same single asymptotically flat end.  The hypersurface is
outer area-minimizing, with disconnected competitors permitted, and
its area $\widehat A=|\Sigma|_{\ghat}$ satisfies
\begin{equation}\label{four:maximal:end:area-comparison}
 \widehat A\ge e^{-3\eps}A_*.
\end{equation}
\end{lemma}

\begin{proof}
For the classical codimension-one regularity background, see
\cite[Sections~5.3--5.4]{Federer1969} and \cite{Federer1970}; the
smooth-Riemannian formulation is
\cite[Chapter~7, \S5, Theorem~5.8, p.~200]{Simon2014}.
Apply Proposition~\ref{four:input:enclosure} to
$(X,g_X)=(\Omega,\ghat)$.  The threshold construction gives a smooth
connected orientable four-dimensional exterior with nonempty compact
smooth boundary $B$ and exactly one end.  Proposition~\ref{four:maximal:prop:end-control}
gives completeness with $B$ included and
$\ghat-\delta=O_2(r^{-q})$ with $q>1$.  These are the full geometric
hypotheses of the cited proposition; it imposes no scalar-curvature
condition and does not require decay exponent two.  Its perimeter counts
the entire frontier, including contact with the obstacle.

The strict inequality $\widehat H_B<0$ places the minimizing frontier
$\Sigma$ strictly outside $B$ and makes it a nonempty compact smooth
embedded minimal hypersurface.  The same proposition gives its connected
complete one-ended exterior and outer area-minimization with every
component counted, including disconnected enclosing competitors.

Every such cut encloses the original obstacle.  The cut comparison in
Lemma~\ref{four:maximal:lem:threshold-collars} gives $|\Sigma|_g\ge A_*$.
Restriction of $\ghat\ge e^{-2\eps}g$ to each three-dimensional tangent
plane gives
\[
 |\Sigma|_{\ghat}\ge e^{-3\eps}|\Sigma|_g
                    \ge e^{-3\eps}A_*.
\]
This proves the asserted area comparison.  The argument uses the
general full-perimeter proposition directly, with its full $q>1$ decay range.  An enclosure result requiring
$O_2(r^{-2})$ decay would not suffice for this application.
\end{proof}

\subsection{The Riemannian inequality and the ordered limits}

We use the following dimension-four instance of the boundary
Riemannian Penrose inequality
\cite[Theorem~1.4, p.~84]{BrayLee2009}.  Let $(\mathscr E,g_{\mathscr E})$ be a smooth
connected four-dimensional Riemannian manifold, complete with its
nonempty compact boundary included, having one asymptotically flat
end.  Suppose in that end
\[
 (g_{\mathscr E})_{ij}-\delta_{ij}=O_2(r^{-p_0}),\quad p_0>1,
 \qquad R_{g_{\mathscr E}}=O(r^{-s_0}),\quad s_0>4,
\]
and $R_{g_{\mathscr E}}\ge0$ everywhere.  If its boundary is minimal and outer
area-minimizing, then
\begin{equation}\label{four:maximal:end:bray-lee}
 E_{g_{\mathscr E}}\ge\frac12
          \left(\frac{|\partial\mathscr E|_{g_{\mathscr E}}}{\omega_3}\right)^{2/3}.
\end{equation}
The boundary may have one or more components, and the ADM
normalization is $1/(6\omega_3)$.  The numerical inequality does not
require a spin assumption; the additional spin hypothesis in the
published theorem concerns its equality conclusion.  We use only
\eqref{four:maximal:end:bray-lee}.

\begin{proof}[Proof of Theorem~\ref{four:maximal:thm:main}]
First keep one set of strict approximation data fixed.  For fixed
$\eps>0$ and each sufficiently large $N$, take the exhaustion limit
of Proposition~\ref{four:maximal:prop:end-control}.  Apply
\eqref{four:maximal:end:bray-lee} to the exterior of $\Sigma$ from
Lemma~\ref{four:maximal:lem:minimal-enclosure}, with metric $\ghat$.
That exterior is smooth and complete with boundary, has one end,
and has nonnegative scalar curvature.  Its decay exponents are
$p_0=q>1$ and $s_0=4+\delta>4$ by \eqref{four:maximal:end:geometry}.  Its boundary
is compact, minimal, and outer area-minimizing by the lemma.  Its
ADM energy exists by Proposition~\ref{four:maximal:prop:flux} and is unchanged
by restriction past the compact region.  Thus every hypothesis of
the Riemannian theorem has been checked for this new exterior.

Equations~\eqref{four:maximal:end:mass-flux}, \eqref{four:maximal:end:flux-deficit}, and
\eqref{four:maximal:end:area-comparison} give
\begin{align*}
 E_g+\frac{\Pi_N}{3\omega_3}
                  \left(\ell+\frac{\ell^2}{\eta_N}\right)
 &\ge E_{\ghat}
 \ge\frac12\left(\frac{\widehat A}{\omega_3}\right)^{2/3}\\
 &\ge\frac12e^{-2\eps}
                       \left(\frac{A_*}{\omega_3}\right)^{2/3}.
\end{align*}
Let $N\to\infty$ with these strict data and $\eps$ fixed.  The error
vanishes by its polynomial bound.  Letting $\eps\downarrow0$ next
yields
\[
 E_g\ge\frac12\left(\frac{A_*}{\omega_3}\right)^{2/3}
\]
for the strict data.  No convergence of the metrics $\ghat$ as
$N\to\infty$ has been assumed: only their individual inequalities
and the uniform deficit are used.

Finally remove the strict approximation using
Proposition~\ref{four:maximal:prop:strict-data}.  The approximation energies
converge to the original $E$, and their enclosing infima $A_*$
converge to the original enclosing infimum.  The given outer
area-minimizing hypothesis, including the admissibility of $S$
itself and all disconnected enclosing competitors, identifies this
infimum with $A=|S|_g$.  Therefore
\[
 \boxed{\displaystyle
 E\ge\frac12\left(\frac{A}{2\pi^2}\right)^{2/3}.}
\]
The prescribed $P=0$ and $E>0$ identify this energy with the ADM rest
mass $m=\sqrt{E^2-|P|^2}=E$.  The original boundary was used with its
future MOTS condition; no minimality or vanishing tangential trace
of $K$ on that boundary was imposed.
\end{proof}
